\documentclass[11pt]{article}

    \usepackage[breakable]{tcolorbox}
    \usepackage{parskip} % Stop auto-indenting (to mimic markdown behaviour)
    

    % Basic figure setup, for now with no caption control since it's done
    % automatically by Pandoc (which extracts ![](path) syntax from Markdown).
    \usepackage{graphicx}
    % Keep aspect ratio if custom image width or height is specified
    \setkeys{Gin}{keepaspectratio}
    % Maintain compatibility with old templates. Remove in nbconvert 6.0
    \let\Oldincludegraphics\includegraphics
    % Ensure that by default, figures have no caption (until we provide a
    % proper Figure object with a Caption API and a way to capture that
    % in the conversion process - todo).
    \usepackage{caption}
    \DeclareCaptionFormat{nocaption}{}
    \captionsetup{format=nocaption,aboveskip=0pt,belowskip=0pt}

    \usepackage{float}
    \floatplacement{figure}{H} % forces figures to be placed at the correct location
    \usepackage{xcolor} % Allow colors to be defined
    \usepackage{enumerate} % Needed for markdown enumerations to work
    \usepackage{geometry} % Used to adjust the document margins
    \usepackage{amsmath} % Equations
    \usepackage{amssymb} % Equations
    \usepackage{textcomp} % defines textquotesingle
    % Hack from http://tex.stackexchange.com/a/47451/13684:
    \AtBeginDocument{%
        \def\PYZsq{\textquotesingle}% Upright quotes in Pygmentized code
    }
    \usepackage{upquote} % Upright quotes for verbatim code
    \usepackage{eurosym} % defines \euro

    \usepackage{iftex}
    \ifPDFTeX
        \usepackage[T1]{fontenc}
        \IfFileExists{alphabeta.sty}{
              \usepackage{alphabeta}
          }{
              \usepackage[mathletters]{ucs}
              \usepackage[utf8x]{inputenc}
          }
    \else
        \usepackage{fontspec}
        \usepackage{unicode-math}
    \fi

    \usepackage{fancyvrb} % verbatim replacement that allows latex
    \usepackage{grffile} % extends the file name processing of package graphics
                         % to support a larger range
    \makeatletter % fix for old versions of grffile with XeLaTeX
    \@ifpackagelater{grffile}{2019/11/01}
    {
      % Do nothing on new versions
    }
    {
      \def\Gread@@xetex#1{%
        \IfFileExists{"\Gin@base".bb}%
        {\Gread@eps{\Gin@base.bb}}%
        {\Gread@@xetex@aux#1}%
      }
    }
    \makeatother
    \usepackage[Export]{adjustbox} % Used to constrain images to a maximum size
    \adjustboxset{max size={0.9\linewidth}{0.9\paperheight}}

    % The hyperref package gives us a pdf with properly built
    % internal navigation ('pdf bookmarks' for the table of contents,
    % internal cross-reference links, web links for URLs, etc.)
    \usepackage{hyperref}
    % The default LaTeX title has an obnoxious amount of whitespace. By default,
    % titling removes some of it. It also provides customization options.
    \usepackage{titling}
    \usepackage{longtable} % longtable support required by pandoc >1.10
    \usepackage{booktabs}  % table support for pandoc > 1.12.2
    \usepackage{array}     % table support for pandoc >= 2.11.3
    \usepackage{calc}      % table minipage width calculation for pandoc >= 2.11.1
    \usepackage[inline]{enumitem} % IRkernel/repr support (it uses the enumerate* environment)
    \usepackage[normalem]{ulem} % ulem is needed to support strikethroughs (\sout)
                                % normalem makes italics be italics, not underlines
    \usepackage{soul}      % strikethrough (\st) support for pandoc >= 3.0.0
    \usepackage{mathrsfs}
    

    
    % Colors for the hyperref package
    \definecolor{urlcolor}{rgb}{0,.145,.698}
    \definecolor{linkcolor}{rgb}{.71,0.21,0.01}
    \definecolor{citecolor}{rgb}{.12,.54,.11}

    % ANSI colors
    \definecolor{ansi-black}{HTML}{3E424D}
    \definecolor{ansi-black-intense}{HTML}{282C36}
    \definecolor{ansi-red}{HTML}{E75C58}
    \definecolor{ansi-red-intense}{HTML}{B22B31}
    \definecolor{ansi-green}{HTML}{00A250}
    \definecolor{ansi-green-intense}{HTML}{007427}
    \definecolor{ansi-yellow}{HTML}{DDB62B}
    \definecolor{ansi-yellow-intense}{HTML}{B27D12}
    \definecolor{ansi-blue}{HTML}{208FFB}
    \definecolor{ansi-blue-intense}{HTML}{0065CA}
    \definecolor{ansi-magenta}{HTML}{D160C4}
    \definecolor{ansi-magenta-intense}{HTML}{A03196}
    \definecolor{ansi-cyan}{HTML}{60C6C8}
    \definecolor{ansi-cyan-intense}{HTML}{258F8F}
    \definecolor{ansi-white}{HTML}{C5C1B4}
    \definecolor{ansi-white-intense}{HTML}{A1A6B2}
    \definecolor{ansi-default-inverse-fg}{HTML}{FFFFFF}
    \definecolor{ansi-default-inverse-bg}{HTML}{000000}

    % common color for the border for error outputs.
    \definecolor{outerrorbackground}{HTML}{FFDFDF}

    % commands and environments needed by pandoc snippets
    % extracted from the output of `pandoc -s`
    \providecommand{\tightlist}{%
      \setlength{\itemsep}{0pt}\setlength{\parskip}{0pt}}
    \DefineVerbatimEnvironment{Highlighting}{Verbatim}{commandchars=\\\{\}}
    % Add ',fontsize=\small' for more characters per line
    \newenvironment{Shaded}{}{}
    \newcommand{\KeywordTok}[1]{\textcolor[rgb]{0.00,0.44,0.13}{\textbf{{#1}}}}
    \newcommand{\DataTypeTok}[1]{\textcolor[rgb]{0.56,0.13,0.00}{{#1}}}
    \newcommand{\DecValTok}[1]{\textcolor[rgb]{0.25,0.63,0.44}{{#1}}}
    \newcommand{\BaseNTok}[1]{\textcolor[rgb]{0.25,0.63,0.44}{{#1}}}
    \newcommand{\FloatTok}[1]{\textcolor[rgb]{0.25,0.63,0.44}{{#1}}}
    \newcommand{\CharTok}[1]{\textcolor[rgb]{0.25,0.44,0.63}{{#1}}}
    \newcommand{\StringTok}[1]{\textcolor[rgb]{0.25,0.44,0.63}{{#1}}}
    \newcommand{\CommentTok}[1]{\textcolor[rgb]{0.38,0.63,0.69}{\textit{{#1}}}}
    \newcommand{\OtherTok}[1]{\textcolor[rgb]{0.00,0.44,0.13}{{#1}}}
    \newcommand{\AlertTok}[1]{\textcolor[rgb]{1.00,0.00,0.00}{\textbf{{#1}}}}
    \newcommand{\FunctionTok}[1]{\textcolor[rgb]{0.02,0.16,0.49}{{#1}}}
    \newcommand{\RegionMarkerTok}[1]{{#1}}
    \newcommand{\ErrorTok}[1]{\textcolor[rgb]{1.00,0.00,0.00}{\textbf{{#1}}}}
    \newcommand{\NormalTok}[1]{{#1}}

    % Additional commands for more recent versions of Pandoc
    \newcommand{\ConstantTok}[1]{\textcolor[rgb]{0.53,0.00,0.00}{{#1}}}
    \newcommand{\SpecialCharTok}[1]{\textcolor[rgb]{0.25,0.44,0.63}{{#1}}}
    \newcommand{\VerbatimStringTok}[1]{\textcolor[rgb]{0.25,0.44,0.63}{{#1}}}
    \newcommand{\SpecialStringTok}[1]{\textcolor[rgb]{0.73,0.40,0.53}{{#1}}}
    \newcommand{\ImportTok}[1]{{#1}}
    \newcommand{\DocumentationTok}[1]{\textcolor[rgb]{0.73,0.13,0.13}{\textit{{#1}}}}
    \newcommand{\AnnotationTok}[1]{\textcolor[rgb]{0.38,0.63,0.69}{\textbf{\textit{{#1}}}}}
    \newcommand{\CommentVarTok}[1]{\textcolor[rgb]{0.38,0.63,0.69}{\textbf{\textit{{#1}}}}}
    \newcommand{\VariableTok}[1]{\textcolor[rgb]{0.10,0.09,0.49}{{#1}}}
    \newcommand{\ControlFlowTok}[1]{\textcolor[rgb]{0.00,0.44,0.13}{\textbf{{#1}}}}
    \newcommand{\OperatorTok}[1]{\textcolor[rgb]{0.40,0.40,0.40}{{#1}}}
    \newcommand{\BuiltInTok}[1]{{#1}}
    \newcommand{\ExtensionTok}[1]{{#1}}
    \newcommand{\PreprocessorTok}[1]{\textcolor[rgb]{0.74,0.48,0.00}{{#1}}}
    \newcommand{\AttributeTok}[1]{\textcolor[rgb]{0.49,0.56,0.16}{{#1}}}
    \newcommand{\InformationTok}[1]{\textcolor[rgb]{0.38,0.63,0.69}{\textbf{\textit{{#1}}}}}
    \newcommand{\WarningTok}[1]{\textcolor[rgb]{0.38,0.63,0.69}{\textbf{\textit{{#1}}}}}
    \makeatletter
    \newsavebox\pandoc@box
    \newcommand*\pandocbounded[1]{%
      \sbox\pandoc@box{#1}%
      % scaling factors for width and height
      \Gscale@div\@tempa\textheight{\dimexpr\ht\pandoc@box+\dp\pandoc@box\relax}%
      \Gscale@div\@tempb\linewidth{\wd\pandoc@box}%
      % select the smaller of both
      \ifdim\@tempb\p@<\@tempa\p@
        \let\@tempa\@tempb
      \fi
      % scaling accordingly (\@tempa < 1)
      \ifdim\@tempa\p@<\p@
        \scalebox{\@tempa}{\usebox\pandoc@box}%
      % scaling not needed, use as it is
      \else
        \usebox{\pandoc@box}%
      \fi
    }
    \makeatother

    % Define a nice break command that doesn't care if a line doesn't already
    % exist.
    \def\br{\hspace*{\fill} \\* }
    % Math Jax compatibility definitions
    \def\gt{>}
    \def\lt{<}
    \let\Oldtex\TeX
    \let\Oldlatex\LaTeX
    \renewcommand{\TeX}{\textrm{\Oldtex}}
    \renewcommand{\LaTeX}{\textrm{\Oldlatex}}
    % Document parameters
    % Document title
    \title{adiabatic}
    
    
    
    
    
    
    
% Pygments definitions
\makeatletter
\def\PY@reset{\let\PY@it=\relax \let\PY@bf=\relax%
    \let\PY@ul=\relax \let\PY@tc=\relax%
    \let\PY@bc=\relax \let\PY@ff=\relax}
\def\PY@tok#1{\csname PY@tok@#1\endcsname}
\def\PY@toks#1+{\ifx\relax#1\empty\else%
    \PY@tok{#1}\expandafter\PY@toks\fi}
\def\PY@do#1{\PY@bc{\PY@tc{\PY@ul{%
    \PY@it{\PY@bf{\PY@ff{#1}}}}}}}
\def\PY#1#2{\PY@reset\PY@toks#1+\relax+\PY@do{#2}}

\@namedef{PY@tok@w}{\def\PY@tc##1{\textcolor[rgb]{0.73,0.73,0.73}{##1}}}
\@namedef{PY@tok@c}{\let\PY@it=\textit\def\PY@tc##1{\textcolor[rgb]{0.24,0.48,0.48}{##1}}}
\@namedef{PY@tok@cp}{\def\PY@tc##1{\textcolor[rgb]{0.61,0.40,0.00}{##1}}}
\@namedef{PY@tok@k}{\let\PY@bf=\textbf\def\PY@tc##1{\textcolor[rgb]{0.00,0.50,0.00}{##1}}}
\@namedef{PY@tok@kp}{\def\PY@tc##1{\textcolor[rgb]{0.00,0.50,0.00}{##1}}}
\@namedef{PY@tok@kt}{\def\PY@tc##1{\textcolor[rgb]{0.69,0.00,0.25}{##1}}}
\@namedef{PY@tok@o}{\def\PY@tc##1{\textcolor[rgb]{0.40,0.40,0.40}{##1}}}
\@namedef{PY@tok@ow}{\let\PY@bf=\textbf\def\PY@tc##1{\textcolor[rgb]{0.67,0.13,1.00}{##1}}}
\@namedef{PY@tok@nb}{\def\PY@tc##1{\textcolor[rgb]{0.00,0.50,0.00}{##1}}}
\@namedef{PY@tok@nf}{\def\PY@tc##1{\textcolor[rgb]{0.00,0.00,1.00}{##1}}}
\@namedef{PY@tok@nc}{\let\PY@bf=\textbf\def\PY@tc##1{\textcolor[rgb]{0.00,0.00,1.00}{##1}}}
\@namedef{PY@tok@nn}{\let\PY@bf=\textbf\def\PY@tc##1{\textcolor[rgb]{0.00,0.00,1.00}{##1}}}
\@namedef{PY@tok@ne}{\let\PY@bf=\textbf\def\PY@tc##1{\textcolor[rgb]{0.80,0.25,0.22}{##1}}}
\@namedef{PY@tok@nv}{\def\PY@tc##1{\textcolor[rgb]{0.10,0.09,0.49}{##1}}}
\@namedef{PY@tok@no}{\def\PY@tc##1{\textcolor[rgb]{0.53,0.00,0.00}{##1}}}
\@namedef{PY@tok@nl}{\def\PY@tc##1{\textcolor[rgb]{0.46,0.46,0.00}{##1}}}
\@namedef{PY@tok@ni}{\let\PY@bf=\textbf\def\PY@tc##1{\textcolor[rgb]{0.44,0.44,0.44}{##1}}}
\@namedef{PY@tok@na}{\def\PY@tc##1{\textcolor[rgb]{0.41,0.47,0.13}{##1}}}
\@namedef{PY@tok@nt}{\let\PY@bf=\textbf\def\PY@tc##1{\textcolor[rgb]{0.00,0.50,0.00}{##1}}}
\@namedef{PY@tok@nd}{\def\PY@tc##1{\textcolor[rgb]{0.67,0.13,1.00}{##1}}}
\@namedef{PY@tok@s}{\def\PY@tc##1{\textcolor[rgb]{0.73,0.13,0.13}{##1}}}
\@namedef{PY@tok@sd}{\let\PY@it=\textit\def\PY@tc##1{\textcolor[rgb]{0.73,0.13,0.13}{##1}}}
\@namedef{PY@tok@si}{\let\PY@bf=\textbf\def\PY@tc##1{\textcolor[rgb]{0.64,0.35,0.47}{##1}}}
\@namedef{PY@tok@se}{\let\PY@bf=\textbf\def\PY@tc##1{\textcolor[rgb]{0.67,0.36,0.12}{##1}}}
\@namedef{PY@tok@sr}{\def\PY@tc##1{\textcolor[rgb]{0.64,0.35,0.47}{##1}}}
\@namedef{PY@tok@ss}{\def\PY@tc##1{\textcolor[rgb]{0.10,0.09,0.49}{##1}}}
\@namedef{PY@tok@sx}{\def\PY@tc##1{\textcolor[rgb]{0.00,0.50,0.00}{##1}}}
\@namedef{PY@tok@m}{\def\PY@tc##1{\textcolor[rgb]{0.40,0.40,0.40}{##1}}}
\@namedef{PY@tok@gh}{\let\PY@bf=\textbf\def\PY@tc##1{\textcolor[rgb]{0.00,0.00,0.50}{##1}}}
\@namedef{PY@tok@gu}{\let\PY@bf=\textbf\def\PY@tc##1{\textcolor[rgb]{0.50,0.00,0.50}{##1}}}
\@namedef{PY@tok@gd}{\def\PY@tc##1{\textcolor[rgb]{0.63,0.00,0.00}{##1}}}
\@namedef{PY@tok@gi}{\def\PY@tc##1{\textcolor[rgb]{0.00,0.52,0.00}{##1}}}
\@namedef{PY@tok@gr}{\def\PY@tc##1{\textcolor[rgb]{0.89,0.00,0.00}{##1}}}
\@namedef{PY@tok@ge}{\let\PY@it=\textit}
\@namedef{PY@tok@gs}{\let\PY@bf=\textbf}
\@namedef{PY@tok@ges}{\let\PY@bf=\textbf\let\PY@it=\textit}
\@namedef{PY@tok@gp}{\let\PY@bf=\textbf\def\PY@tc##1{\textcolor[rgb]{0.00,0.00,0.50}{##1}}}
\@namedef{PY@tok@go}{\def\PY@tc##1{\textcolor[rgb]{0.44,0.44,0.44}{##1}}}
\@namedef{PY@tok@gt}{\def\PY@tc##1{\textcolor[rgb]{0.00,0.27,0.87}{##1}}}
\@namedef{PY@tok@err}{\def\PY@bc##1{{\setlength{\fboxsep}{\string -\fboxrule}\fcolorbox[rgb]{1.00,0.00,0.00}{1,1,1}{\strut ##1}}}}
\@namedef{PY@tok@kc}{\let\PY@bf=\textbf\def\PY@tc##1{\textcolor[rgb]{0.00,0.50,0.00}{##1}}}
\@namedef{PY@tok@kd}{\let\PY@bf=\textbf\def\PY@tc##1{\textcolor[rgb]{0.00,0.50,0.00}{##1}}}
\@namedef{PY@tok@kn}{\let\PY@bf=\textbf\def\PY@tc##1{\textcolor[rgb]{0.00,0.50,0.00}{##1}}}
\@namedef{PY@tok@kr}{\let\PY@bf=\textbf\def\PY@tc##1{\textcolor[rgb]{0.00,0.50,0.00}{##1}}}
\@namedef{PY@tok@bp}{\def\PY@tc##1{\textcolor[rgb]{0.00,0.50,0.00}{##1}}}
\@namedef{PY@tok@fm}{\def\PY@tc##1{\textcolor[rgb]{0.00,0.00,1.00}{##1}}}
\@namedef{PY@tok@vc}{\def\PY@tc##1{\textcolor[rgb]{0.10,0.09,0.49}{##1}}}
\@namedef{PY@tok@vg}{\def\PY@tc##1{\textcolor[rgb]{0.10,0.09,0.49}{##1}}}
\@namedef{PY@tok@vi}{\def\PY@tc##1{\textcolor[rgb]{0.10,0.09,0.49}{##1}}}
\@namedef{PY@tok@vm}{\def\PY@tc##1{\textcolor[rgb]{0.10,0.09,0.49}{##1}}}
\@namedef{PY@tok@sa}{\def\PY@tc##1{\textcolor[rgb]{0.73,0.13,0.13}{##1}}}
\@namedef{PY@tok@sb}{\def\PY@tc##1{\textcolor[rgb]{0.73,0.13,0.13}{##1}}}
\@namedef{PY@tok@sc}{\def\PY@tc##1{\textcolor[rgb]{0.73,0.13,0.13}{##1}}}
\@namedef{PY@tok@dl}{\def\PY@tc##1{\textcolor[rgb]{0.73,0.13,0.13}{##1}}}
\@namedef{PY@tok@s2}{\def\PY@tc##1{\textcolor[rgb]{0.73,0.13,0.13}{##1}}}
\@namedef{PY@tok@sh}{\def\PY@tc##1{\textcolor[rgb]{0.73,0.13,0.13}{##1}}}
\@namedef{PY@tok@s1}{\def\PY@tc##1{\textcolor[rgb]{0.73,0.13,0.13}{##1}}}
\@namedef{PY@tok@mb}{\def\PY@tc##1{\textcolor[rgb]{0.40,0.40,0.40}{##1}}}
\@namedef{PY@tok@mf}{\def\PY@tc##1{\textcolor[rgb]{0.40,0.40,0.40}{##1}}}
\@namedef{PY@tok@mh}{\def\PY@tc##1{\textcolor[rgb]{0.40,0.40,0.40}{##1}}}
\@namedef{PY@tok@mi}{\def\PY@tc##1{\textcolor[rgb]{0.40,0.40,0.40}{##1}}}
\@namedef{PY@tok@il}{\def\PY@tc##1{\textcolor[rgb]{0.40,0.40,0.40}{##1}}}
\@namedef{PY@tok@mo}{\def\PY@tc##1{\textcolor[rgb]{0.40,0.40,0.40}{##1}}}
\@namedef{PY@tok@ch}{\let\PY@it=\textit\def\PY@tc##1{\textcolor[rgb]{0.24,0.48,0.48}{##1}}}
\@namedef{PY@tok@cm}{\let\PY@it=\textit\def\PY@tc##1{\textcolor[rgb]{0.24,0.48,0.48}{##1}}}
\@namedef{PY@tok@cpf}{\let\PY@it=\textit\def\PY@tc##1{\textcolor[rgb]{0.24,0.48,0.48}{##1}}}
\@namedef{PY@tok@c1}{\let\PY@it=\textit\def\PY@tc##1{\textcolor[rgb]{0.24,0.48,0.48}{##1}}}
\@namedef{PY@tok@cs}{\let\PY@it=\textit\def\PY@tc##1{\textcolor[rgb]{0.24,0.48,0.48}{##1}}}

\def\PYZbs{\char`\\}
\def\PYZus{\char`\_}
\def\PYZob{\char`\{}
\def\PYZcb{\char`\}}
\def\PYZca{\char`\^}
\def\PYZam{\char`\&}
\def\PYZlt{\char`\<}
\def\PYZgt{\char`\>}
\def\PYZsh{\char`\#}
\def\PYZpc{\char`\%}
\def\PYZdl{\char`\$}
\def\PYZhy{\char`\-}
\def\PYZsq{\char`\'}
\def\PYZdq{\char`\"}
\def\PYZti{\char`\~}
% for compatibility with earlier versions
\def\PYZat{@}
\def\PYZlb{[}
\def\PYZrb{]}
\makeatother


    % For linebreaks inside Verbatim environment from package fancyvrb.
    \makeatletter
        \newbox\Wrappedcontinuationbox
        \newbox\Wrappedvisiblespacebox
        \newcommand*\Wrappedvisiblespace {\textcolor{red}{\textvisiblespace}}
        \newcommand*\Wrappedcontinuationsymbol {\textcolor{red}{\llap{\tiny$\m@th\hookrightarrow$}}}
        \newcommand*\Wrappedcontinuationindent {3ex }
        \newcommand*\Wrappedafterbreak {\kern\Wrappedcontinuationindent\copy\Wrappedcontinuationbox}
        % Take advantage of the already applied Pygments mark-up to insert
        % potential linebreaks for TeX processing.
        %        {, <, #, %, $, ' and ": go to next line.
        %        _, }, ^, &, >, - and ~: stay at end of broken line.
        % Use of \textquotesingle for straight quote.
        \newcommand*\Wrappedbreaksatspecials {%
            \def\PYGZus{\discretionary{\char`\_}{\Wrappedafterbreak}{\char`\_}}%
            \def\PYGZob{\discretionary{}{\Wrappedafterbreak\char`\{}{\char`\{}}%
            \def\PYGZcb{\discretionary{\char`\}}{\Wrappedafterbreak}{\char`\}}}%
            \def\PYGZca{\discretionary{\char`\^}{\Wrappedafterbreak}{\char`\^}}%
            \def\PYGZam{\discretionary{\char`\&}{\Wrappedafterbreak}{\char`\&}}%
            \def\PYGZlt{\discretionary{}{\Wrappedafterbreak\char`\<}{\char`\<}}%
            \def\PYGZgt{\discretionary{\char`\>}{\Wrappedafterbreak}{\char`\>}}%
            \def\PYGZsh{\discretionary{}{\Wrappedafterbreak\char`\#}{\char`\#}}%
            \def\PYGZpc{\discretionary{}{\Wrappedafterbreak\char`\%}{\char`\%}}%
            \def\PYGZdl{\discretionary{}{\Wrappedafterbreak\char`\$}{\char`\$}}%
            \def\PYGZhy{\discretionary{\char`\-}{\Wrappedafterbreak}{\char`\-}}%
            \def\PYGZsq{\discretionary{}{\Wrappedafterbreak\textquotesingle}{\textquotesingle}}%
            \def\PYGZdq{\discretionary{}{\Wrappedafterbreak\char`\"}{\char`\"}}%
            \def\PYGZti{\discretionary{\char`\~}{\Wrappedafterbreak}{\char`\~}}%
        }
        % Some characters . , ; ? ! / are not pygmentized.
        % This macro makes them "active" and they will insert potential linebreaks
        \newcommand*\Wrappedbreaksatpunct {%
            \lccode`\~`\.\lowercase{\def~}{\discretionary{\hbox{\char`\.}}{\Wrappedafterbreak}{\hbox{\char`\.}}}%
            \lccode`\~`\,\lowercase{\def~}{\discretionary{\hbox{\char`\,}}{\Wrappedafterbreak}{\hbox{\char`\,}}}%
            \lccode`\~`\;\lowercase{\def~}{\discretionary{\hbox{\char`\;}}{\Wrappedafterbreak}{\hbox{\char`\;}}}%
            \lccode`\~`\:\lowercase{\def~}{\discretionary{\hbox{\char`\:}}{\Wrappedafterbreak}{\hbox{\char`\:}}}%
            \lccode`\~`\?\lowercase{\def~}{\discretionary{\hbox{\char`\?}}{\Wrappedafterbreak}{\hbox{\char`\?}}}%
            \lccode`\~`\!\lowercase{\def~}{\discretionary{\hbox{\char`\!}}{\Wrappedafterbreak}{\hbox{\char`\!}}}%
            \lccode`\~`\/\lowercase{\def~}{\discretionary{\hbox{\char`\/}}{\Wrappedafterbreak}{\hbox{\char`\/}}}%
            \catcode`\.\active
            \catcode`\,\active
            \catcode`\;\active
            \catcode`\:\active
            \catcode`\?\active
            \catcode`\!\active
            \catcode`\/\active
            \lccode`\~`\~
        }
    \makeatother

    \let\OriginalVerbatim=\Verbatim
    \makeatletter
    \renewcommand{\Verbatim}[1][1]{%
        %\parskip\z@skip
        \sbox\Wrappedcontinuationbox {\Wrappedcontinuationsymbol}%
        \sbox\Wrappedvisiblespacebox {\FV@SetupFont\Wrappedvisiblespace}%
        \def\FancyVerbFormatLine ##1{\hsize\linewidth
            \vtop{\raggedright\hyphenpenalty\z@\exhyphenpenalty\z@
                \doublehyphendemerits\z@\finalhyphendemerits\z@
                \strut ##1\strut}%
        }%
        % If the linebreak is at a space, the latter will be displayed as visible
        % space at end of first line, and a continuation symbol starts next line.
        % Stretch/shrink are however usually zero for typewriter font.
        \def\FV@Space {%
            \nobreak\hskip\z@ plus\fontdimen3\font minus\fontdimen4\font
            \discretionary{\copy\Wrappedvisiblespacebox}{\Wrappedafterbreak}
            {\kern\fontdimen2\font}%
        }%

        % Allow breaks at special characters using \PYG... macros.
        \Wrappedbreaksatspecials
        % Breaks at punctuation characters . , ; ? ! and / need catcode=\active
        \OriginalVerbatim[#1,codes*=\Wrappedbreaksatpunct]%
    }
    \makeatother

    % Exact colors from NB
    \definecolor{incolor}{HTML}{303F9F}
    \definecolor{outcolor}{HTML}{D84315}
    \definecolor{cellborder}{HTML}{CFCFCF}
    \definecolor{cellbackground}{HTML}{F7F7F7}

    % prompt
    \makeatletter
    \newcommand{\boxspacing}{\kern\kvtcb@left@rule\kern\kvtcb@boxsep}
    \makeatother
    \newcommand{\prompt}[4]{
        {\ttfamily\llap{{\color{#2}[#3]:\hspace{3pt}#4}}\vspace{-\baselineskip}}
    }
    

    

    % Prevent overflowing lines due to hard-to-break entities
    \sloppy
    % Setup hyperref package
    \hypersetup{
      breaklinks=true,  % so long urls are correctly broken across lines
      colorlinks=true,
      urlcolor=urlcolor,
      linkcolor=linkcolor,
      citecolor=citecolor,
      }
    % Slightly bigger margins than the latex defaults
    
    \geometry{verbose,tmargin=1in,bmargin=1in,lmargin=1in,rmargin=1in}
    
    

% Quantum bra/ket macros (added by qs-latex template)
\providecommand{\ket}[1]{\left|#1\right\rangle}
\providecommand{\bra}[1]{\left\langle#1\right|}
\providecommand{\braket}[2]{\left\langle#1\middle|#2\right\rangle}


\begin{document}
    
    \maketitle
    
    

    
    \section{Ising Adiabatic State
Preparation}\label{ising-adiabatic-state-preparation}

\begin{center}\rule{0.5\linewidth}{0.5pt}\end{center}

\subsection{\texorpdfstring{\emph{Review}: Postulate
6}{Review: Postulate 6}}\label{review-postulate-6}

\begin{quote}
The time evolution of the state vector \(\ket{\psi(t)}\) is governed by
the Schrödinger equation:
\[      i\hbar\frac{d}{dt}\ket{\psi(t)} = \hat{H}(t)\ket{\psi(t)},\]
\end{quote}

\subsubsection{Time independent case}\label{time-independent-case}

Imagine the case where \(\hat{H}(t) \neq f(t)\). This is just a
statement that the Hamiltoninan operator doesn't change in time. If that
is the case, then we can directly integrate this differential equation,
to get the following solution: \[\begin{align}
\ket{\psi(t)} =& e^{\frac{-i}{\hbar}\hat{H}\Delta t}\ket{\psi(t_0)}, \\
=& U(t,t_0)\ket{\psi(t_0)}
\end{align}
\] where \(\Delta t = t-t_0\). Differentiate this wavefunction to
convince yourself that this is a solution to the time dependent
Schrödinger equation.

Because of this simple form, we can do fancy things, like first go to
\(t_1\), then go to \(t_2\), e.g., \[\begin{align}
\ket{\psi(t_2)} = U(t_2,t_1)U(t_1,t_0)\ket{\psi(t_0)}
\end{align}
\]

\subsubsection{Time dependent case}\label{time-dependent-case}

The more general (and complicated) case where \(\hat{H}\) is
time-dependent cannot be written down in such a simple form. However, we
can develop an arbitrarily accurate approximation to the solution by
using the result above. Assume we want to solve for \(\ket{\psi(t)}\),
we can break up the evolution into \emph{infinitely} many steps:
\[\begin{align}
\ket{\psi(t)} = \lim_{N\rightarrow \infty} U(t,t_{N-1})\dots U(t_2,t_1) U(t_1,t_0)\ket{\psi(t_0)}
\end{align}\] By dividing the evolution into small enough time steps, we
can use the time independent solution! The core idea, is that within a
small enough time window, \(t_i \rightarrow t_{i+1}\), the Hamiltonian
is changing slowly enough such that it appears constant. This means that
whatever the Hamiltonian looks like at a give time, (\(\hat{H}(t_i)\)),
we can exponentiate it, to obtain a unitary \(U(t_{i+1}, t_i)\).

In practice, of course, we can never write down a real \(\infty\) of
terms. And so we simply stop at some point and accept that there is some
approximation being made. \[\begin{align}
\ket{\psi(t)} \approx   U(t,t_{N-1})\dots U(t_2,t_1) U(t_1,t_0)\ket{\psi(t_0)}
\end{align}\]

We often refer to this particular approximate form of the time evolution
as being \texttt{trotterized}

\begin{center}\rule{0.5\linewidth}{0.5pt}\end{center}

    \section{Adiabatic Principle}\label{adiabatic-principle}

We can use the above result (the ``trotterized'' form of the TDSE
solution) to solve a generally difficult problem, \textbf{finding the
ground state of a Hamiltonian}. In our Ising case, this amounts to
finding the lowest energy spin configuration out of an exponentially
large number of candidates. However, if we were to be able to find the
ground state, we would simultaneously be solving a wide range of
extremely important problems including the travelling salesman problem
or finding the max-cut of a graph or network.

In this notebook, we will use the \texttt{adiabatic\ principle} to
define an interesting (yet simple) quantum algorithm for obtaining
arbitrarily accurate approximations to Hamiltonian ground states.

We will define the adiabatic principle as follows: \textgreater If we
start out in the ground state of a Hamiltonian \(\hat{H}_A\), and the
\emph{slowly} move to Hamiltonian \(\hat{H}_B\), then we will end up in
the ground state of \(\hat{H}_B\) (assuming we move slow enough and
there is a gap).

To use the adiabatic principle to obtain the ground state for a target
Hamiltonian, \(\hat{H}_B\), (which could be the Ising example we have
been working with or some other problem) we will define the following
time dependent Hamiltonian: \[\begin{align}
\hat{H}(t) = (1-f(t)) \hat{H}_A + f(t) \hat{H}_B,
\end{align}\] where \(\hat{H}_A\) is our starting Hamiltonian (this
should be an easy Hamiltonian for which we know the ground state).
Notice that while both \(\hat{H}_A\) and \(\hat{H}_B\) are
time-\emph{independent}, the coefficent \(f(t)\) changes in time, so the
total Hamiltonian is time-\emph{dependent}.

This means that we can define any function \(f(t)\) that starts at 0 and
ends up equal to 1 at some final time, \(t_f\). - \(f(t=0) = 0\) -
\(f(t=T) = 1\)

The simplest time-dependent function that we might consider is jut a
linear interpolation between the two Hamiltonians:

\begin{itemize}
\tightlist
\item
  \$ f(t) = t/T\$ where \(T\) is the total time of evolution. This means
  that a larger \(T\) results in a \emph{slower} evolution, and
  satisfies the condition that \(\hat{H}(0)=\hat{H}_A\) and
  \(\hat{H}(T)=\hat{H}_B\).
\end{itemize}

\subsection{Trotterization}\label{trotterization}

In addition to requiring that our final time \(T\) is large enough to
ensure that our process is slow, we also need to make sure that we are
solving the resulting time dependent schr"odinger equation accurately
enough. Since we are using the trotter approximation (meaning treating
the total evolution as a sequence of short time-independent evolutions),
we will also need to choose a number of time steps \(N\) to be large
enough for the adiabatic evolution to be accurate.

Each time step is evolved by a time step given by: - \(\Delta t = T/N\)

We will want both \(T\) and \(N\) to be large enough for accurate
results.

\[
\begin{align}
\ket{\psi(T)} &\approx U((N+1)\Delta t, N\Delta t)\cdots U(2\Delta t, \Delta t)U(\Delta t, 0)\ket{\psi(0)} \\
\ket{\psi(T)} &\approx \prod_{k=0}^{N-1} U((k+1)\Delta t, k\Delta t)\ket{\psi(0)} = \prod_{k=0}^{N-1}U_k\ket{\psi(0)}\\
\end{align}
\]

\[
\begin{align}
U_k &= e^{-i\Delta t \hat{H}(k\Delta t)}
% U(t+\Delta t, t) &\approx e^{-i\Delta t \left(1-t/T\right)\hat{H}_A} e^{-i\Delta t \left(t/T-\Delta t/2T\right)\hat{H}_B /2}
\end{align}
\]

Since we have an explicit form of our time dependent hamiltonian
\(\hat{H}(t)\), we can write this out: \[
\begin{align}
\hat{H}(k\Delta t) =& (1-\frac{k\Delta t}{T})\hat{H}_A + \frac{k\Delta t}{T}\hat{H}_B
\end{align}
\]

Convince yourself that when \(k=0\), we have the initial hamiltonian,
\(H_A\), and when \(k=N\), we have the final Hamiltonian, \(H_B\)

\subsection{One more trotterization}\label{one-more-trotterization}

\[
\begin{align}
U_k &= e^{-i\Delta t \hat{H}(k\Delta t)} \\
&= e^{-i\Delta t \left( (1-\frac{k\Delta t}{T})\hat{H}_A + \frac{k\Delta t}{T}\hat{H}_B\right)} \\
&\approx e^{-i\Delta t (1-\frac{k\Delta t}{T})\hat{H}_A}  e^{-i\Delta t \frac{k\Delta t}{T}\hat{H}_B} 
\end{align}
\]

    Using \(T=2\) and \(N=4\), (such that \(\Delta t = \tfrac{1}{2}\)) from
the case above we can define the solution as: \[
\begin{align}
\ket{\psi(2)} &= 
e^{-i\Delta t\left(\tfrac{1}{4}\hat{H}_A+\tfrac{3}{4}\hat{H}_B\right)}
e^{-i\Delta t\left(\tfrac{1}{2}\hat{H}_A+\tfrac{1}{2}\hat{H}_B\right)}
e^{-i\Delta t\left(\tfrac{3}{4}\hat{H}_A+\tfrac{1}{4}\hat{H}_B\right)}
e^{-i\Delta t\hat{H}_A} 
\ket{\psi(0)}\\\nonumber \\
&= 
U(2,\tfrac{3}{2})
U(\tfrac{3}{2},1)
U(1,\tfrac{1}{2})
U(\tfrac{1}{2},0) \\\nonumber \\
&\approx 
e^{-i\Delta t\tfrac{1}{4}\hat{H}_A}
e^{-i\Delta t\tfrac{3}{4}\hat{H}_B}
e^{-i\Delta t\tfrac{1}{2}\hat{H}_A}
e^{-i\Delta t\tfrac{1}{2}\hat{H}_B}
e^{-i\Delta t\tfrac{3}{4}\hat{H}_A}
e^{-i\Delta t\tfrac{1}{4}\hat{H}_B}
e^{-i\Delta t\hat{H}_A}
\ket{\psi(0)}\\
\end{align}
\]

    We can use a quantum computer to carry out these operations. But first
we will need to know how to represent these quantities on a quantum
computer.

\section{Quantum Representation of
Ising}\label{quantum-representation-of-ising}

In the previous examples, we have used `classical' bitstrings, by which
I mean all of our states are always written as a single product of zeros
and ones, i.e., \(\ket{010010}\). So far, we have used the Ising
Hamiltonian,
\[ \hat{H} = \sum_{(i,j)\in E}J_{ij} s_is_j + \sum_i \mu_i s_i,\] to
compute the energy of these bitstrings where the spin function takes in
a bit and returns either 1 or -1. For example, if
\(\ket{\psi}=\ket{0101}\), then \(s_2\ket{0101}=0\).

We can make this a bit more mathematically useful, and define our
quantities in terms of linear algebra, making our spin function and
operator, and the state a vector. Each bit is two-dimensional (it can be
either zero or one), so we can express it as a 2-dimensional vector:

\[
\ket{0} = \begin{pmatrix}1 \\ 0\end{pmatrix},
\text{ }~ 
\ket{1} = \begin{pmatrix}0 \\ 1\end{pmatrix} 
\]

Similarly, our ``spin-functions'', \(s_i\), can also be expressed in
terms of a simple, linear algebra quantity, \(\hat{z}\)\\
\[\begin{align}
\hat{z} = \begin{pmatrix}1 & 0 \\ 0 &-1\end{pmatrix}
\end{align}
\]

Consider the following operator, \(\hat{x}\): \[\begin{align}
\hat{x} = \begin{pmatrix}0 & 1 \\ 1 &0\end{pmatrix}
\end{align}
\]

What happens when we apply \(\hat{x}\) to our states? \[
\begin{align}
\hat{x}\ket{0} =& \ket{1} \\ 
\hat{x}\ket{1} =& \ket{0} 
\end{align}
\]

Here, \(\hat{z}\) tells us about the spin being spin-up or spin down.
And \(\hat{x}\), flips the state between up and down. However, because
our bits are quantum states (vectors), they can also point in other
directions besides up and down - this is a key difference between
quantum an classical computers!

Consider the last operator/gate we will need, \(\hat{U}^H\) (the
Hadamard gate): \[
\begin{align}
\hat{U}^H =& \frac{1}{\sqrt{2}}\begin{pmatrix}1 & 1 \\ 1 &-1\end{pmatrix} \\
    =& \frac{1}{\sqrt{2}}\left(\hat{x}+\hat{z}\right)
\end{align}
\]

Applying this operator to a state does something interesting: \[
\begin{align}
\hat{U}^H\ket{0} =& \frac{1}{\sqrt{2}}\left(\hat{x}+\hat{z}\right)\ket{0} \\
=& \frac{1}{\sqrt{2}}\left(\ket{1} + \ket{0}\right) \\  
\hat{U}^H\ket{1} =& \frac{1}{\sqrt{2}}\left(\hat{x}+\hat{z}\right)\ket{1} \\
=& \frac{1}{\sqrt{2}}\left(\ket{0} - \ket{1}\right) \\  
\end{align}
\]

We've now created \texttt{superposition} states! Play around with these
objects in the following python cell.

    \begin{tcolorbox}[breakable, size=fbox, boxrule=1pt, pad at break*=1mm,colback=cellbackground, colframe=cellborder]
\prompt{In}{incolor}{1}{\boxspacing}
\begin{Verbatim}[commandchars=\\\{\}]
\PY{k+kn}{import}\PY{+w}{ }\PY{n+nn}{numpy}\PY{+w}{ }\PY{k}{as}\PY{+w}{ }\PY{n+nn}{np}

\PY{c+c1}{\PYZsh{} define our qubit states}
\PY{n}{s0} \PY{o}{=} \PY{n}{np}\PY{o}{.}\PY{n}{array}\PY{p}{(}\PY{p}{[}\PY{p}{[}\PY{l+m+mi}{1}\PY{p}{]}\PY{p}{,}\PY{p}{[}\PY{l+m+mi}{0}\PY{p}{]}\PY{p}{]}\PY{p}{)}
\PY{n}{s1} \PY{o}{=} \PY{n}{np}\PY{o}{.}\PY{n}{array}\PY{p}{(}\PY{p}{[}\PY{p}{[}\PY{l+m+mi}{0}\PY{p}{]}\PY{p}{,}\PY{p}{[}\PY{l+m+mi}{1}\PY{p}{]}\PY{p}{]}\PY{p}{)}

\PY{c+c1}{\PYZsh{} define z     }
\PY{n}{z} \PY{o}{=} \PY{n}{np}\PY{o}{.}\PY{n}{array}\PY{p}{(}\PY{p}{[}\PY{p}{[}\PY{l+m+mi}{1}\PY{p}{,}\PY{l+m+mi}{0}\PY{p}{]}\PY{p}{,}\PY{p}{[}\PY{l+m+mi}{0}\PY{p}{,}\PY{o}{\PYZhy{}}\PY{l+m+mi}{1}\PY{p}{]}\PY{p}{]}\PY{p}{)}
\PY{n+nb}{print}\PY{p}{(}\PY{l+s+s2}{\PYZdq{}}\PY{l+s+s2}{ z operator: }\PY{l+s+se}{\PYZbs{}n}\PY{l+s+s2}{\PYZdq{}}\PY{p}{,} \PY{n}{z}\PY{p}{)}

\PY{c+c1}{\PYZsh{} define x}
\PY{n}{x} \PY{o}{=} \PY{n}{np}\PY{o}{.}\PY{n}{array}\PY{p}{(}\PY{p}{[}\PY{p}{[}\PY{l+m+mi}{0}\PY{p}{,}\PY{l+m+mi}{1}\PY{p}{]}\PY{p}{,}\PY{p}{[}\PY{l+m+mi}{1}\PY{p}{,}\PY{l+m+mi}{0}\PY{p}{]}\PY{p}{]}\PY{p}{)}
\PY{n+nb}{print}\PY{p}{(}\PY{l+s+s2}{\PYZdq{}}\PY{l+s+s2}{ x operator: }\PY{l+s+se}{\PYZbs{}n}\PY{l+s+s2}{\PYZdq{}}\PY{p}{,} \PY{n}{x}\PY{p}{)}

\PY{c+c1}{\PYZsh{} define Uh }
\PY{n}{Uh} \PY{o}{=} \PY{n}{np}\PY{o}{.}\PY{n}{array}\PY{p}{(}\PY{p}{[}\PY{p}{[}\PY{l+m+mi}{1}\PY{p}{,}\PY{l+m+mi}{1}\PY{p}{]}\PY{p}{,}\PY{p}{[}\PY{l+m+mi}{1}\PY{p}{,}\PY{o}{\PYZhy{}}\PY{l+m+mi}{1}\PY{p}{]}\PY{p}{]}\PY{p}{)}\PY{o}{/}\PY{n}{np}\PY{o}{.}\PY{n}{sqrt}\PY{p}{(}\PY{l+m+mi}{2}\PY{p}{)}
\PY{n+nb}{print}\PY{p}{(}\PY{l+s+s2}{\PYZdq{}}\PY{l+s+s2}{ Uh operator: }\PY{l+s+se}{\PYZbs{}n}\PY{l+s+s2}{\PYZdq{}}\PY{p}{,} \PY{n}{Uh}\PY{p}{)}

\PY{n+nb}{print}\PY{p}{(}\PY{l+s+s2}{\PYZdq{}}\PY{l+s+se}{\PYZbs{}n}\PY{l+s+s2}{ z|0\PYZgt{}: }\PY{l+s+se}{\PYZbs{}n}\PY{l+s+s2}{\PYZdq{}}\PY{p}{,} \PY{n}{z}\PY{n+nd}{@s0}\PY{p}{)}
\PY{n+nb}{print}\PY{p}{(}\PY{l+s+s2}{\PYZdq{}}\PY{l+s+se}{\PYZbs{}n}\PY{l+s+s2}{ z|1\PYZgt{}: }\PY{l+s+se}{\PYZbs{}n}\PY{l+s+s2}{\PYZdq{}}\PY{p}{,} \PY{n}{z}\PY{n+nd}{@s1}\PY{p}{)}

\PY{n+nb}{print}\PY{p}{(}\PY{l+s+s2}{\PYZdq{}}\PY{l+s+se}{\PYZbs{}n}\PY{l+s+s2}{ x|0\PYZgt{}: }\PY{l+s+se}{\PYZbs{}n}\PY{l+s+s2}{\PYZdq{}}\PY{p}{,} \PY{n}{x}\PY{n+nd}{@s0}\PY{p}{)}
\PY{n+nb}{print}\PY{p}{(}\PY{l+s+s2}{\PYZdq{}}\PY{l+s+se}{\PYZbs{}n}\PY{l+s+s2}{ x|1\PYZgt{}: }\PY{l+s+se}{\PYZbs{}n}\PY{l+s+s2}{\PYZdq{}}\PY{p}{,} \PY{n}{x}\PY{n+nd}{@s1}\PY{p}{)}

\PY{n+nb}{print}\PY{p}{(}\PY{l+s+s2}{\PYZdq{}}\PY{l+s+se}{\PYZbs{}n}\PY{l+s+s2}{ Uh|0\PYZgt{}: }\PY{l+s+se}{\PYZbs{}n}\PY{l+s+s2}{\PYZdq{}}\PY{p}{,} \PY{n}{Uh}\PY{n+nd}{@s0}\PY{p}{)}
\PY{n+nb}{print}\PY{p}{(}\PY{l+s+s2}{\PYZdq{}}\PY{l+s+se}{\PYZbs{}n}\PY{l+s+s2}{ Uh|1\PYZgt{}: }\PY{l+s+se}{\PYZbs{}n}\PY{l+s+s2}{\PYZdq{}}\PY{p}{,} \PY{n}{Uh}\PY{n+nd}{@s1}\PY{p}{)}

\PY{n+nb}{print}\PY{p}{(}\PY{l+s+s2}{\PYZdq{}}\PY{l+s+se}{\PYZbs{}n}\PY{l+s+s2}{ |0\PYZgt{}⊗|1\PYZgt{}: }\PY{l+s+se}{\PYZbs{}n}\PY{l+s+s2}{\PYZdq{}}\PY{p}{,} \PY{n}{np}\PY{o}{.}\PY{n}{kron}\PY{p}{(}\PY{n}{s0}\PY{p}{,}\PY{n}{s1}\PY{p}{)}\PY{p}{)}
\PY{n+nb}{print}\PY{p}{(}\PY{l+s+s2}{\PYZdq{}}\PY{l+s+se}{\PYZbs{}n}\PY{l+s+s2}{ z⊗z    : }\PY{l+s+se}{\PYZbs{}n}\PY{l+s+s2}{\PYZdq{}}\PY{p}{,} \PY{n}{np}\PY{o}{.}\PY{n}{kron}\PY{p}{(}\PY{n}{z}\PY{p}{,}\PY{n}{z}\PY{p}{)}\PY{p}{)}
\end{Verbatim}
\end{tcolorbox}

    \begin{Verbatim}[commandchars=\\\{\}]
 z operator:
 [[ 1  0]
 [ 0 -1]]
 x operator:
 [[0 1]
 [1 0]]
 Uh operator:
 [[ 0.70710678  0.70710678]
 [ 0.70710678 -0.70710678]]

 z|0>:
 [[1]
 [0]]

 z|1>:
 [[ 0]
 [-1]]

 x|0>:
 [[0]
 [1]]

 x|1>:
 [[1]
 [0]]

 Uh|0>:
 [[0.70710678]
 [0.70710678]]

 Uh|1>:
 [[ 0.70710678]
 [-0.70710678]]

 |0>⊗|1>:
 [[0]
 [1]
 [0]
 [0]]

 z⊗z    :
 [[ 1  0  0  0]
 [ 0 -1  0  0]
 [ 0  0 -1  0]
 [ 0  0  0  1]]
    \end{Verbatim}

    \section{\texorpdfstring{\textbf{Question}:}{Question:}}\label{question}

In the next cell: explain what \(\hat{U}^H\) does to the states,
\(\ket{0}\) and \(\ket{1}\) in terms of \(\hat{x}\).:

    \(\hat{U}^H\) maps the computational-basis states (\(\hat{z}\)
eigenstates) to the eigenstates of \(\hat{x}\):

\begin{itemize}
\tightlist
\item
  \(\hat{U}^H\ket{0} = \tfrac{1}{\sqrt{2}}(\ket{0}+\ket{1}) \equiv \ket{+}\),
  which is the 1 eigenstate of \(\hat{x}\):
  \(\hat{x}\ket{+} = +\ket{+}\).
\item
  \(\hat{U}^H\ket{1} = \tfrac{1}{\sqrt{2}}(\ket{0}-\ket{1}) \equiv \ket{-}\),
  which is the −1 eigenstate of \(\hat{x}\):
  \(\hat{x}\ket{-} = -\ket{-}\).
\end{itemize}

So \(\hat{U}^H\) rotates the \(\hat{z}\)-basis into the
\(\hat{x}\)-basis. Equivalently,
\(\hat{U}^H \hat{x} \hat{U}^H = \hat{z}\), so applying \(\hat{U}^H\)
before a \(\hat{z}\)-basis measurement is the same as measuring
\(\hat{x}\).

    \begin{center}\rule{0.5\linewidth}{0.5pt}\end{center}

\subsection{Adiabatic Evolution of Ising Hamiltonian with Quantum
Circuit}\label{adiabatic-evolution-of-ising-hamiltonian-with-quantum-circuit}

We discussed that the adiabatic principle can be used to obtain the
ground state of a difficult problem (e.g., Ising Hamiltonian) by
starting with an easy problem (which for which we know the solution) and
slowly moving toward the difficult solution. This means that we need to
\emph{drive our system forward in time} with a time-dependent
Hamiltonian. This is generally a difficult task, but we can leverage a
quantum computer to do this for us!

By choosing \(\hat{H}_0 = - \sum_i \sigma_i^x\), recall that the ground
state of the Ising model can be prepared by the following sequence of
time evolutions:

\[
\ket{\psi(T)} \approx 
U(4\Delta t,3 \Delta t)
U(3\Delta t,2 \Delta t)
U(2\Delta t,\Delta t)
U(\Delta t,0)
\ket{\psi(0)}
\]

which is, \[
\ket{\psi(T)} \approx 
e^{-i\left(\tfrac{1}{4}\hat{H}_0+\tfrac{3}{4}\hat{H}_1\right)\Delta t}
e^{-i\left(\tfrac{1}{2}\hat{H}_0+\tfrac{1}{2}\hat{H}_1\right)\Delta t}
e^{-i\left(\tfrac{3}{4}\hat{H}_0+\tfrac{1}{4}\hat{H}_1\right)\Delta t}
e^{-i\hat{H}_0\Delta t}
\ket{\psi(0)},
\]

where \(N\Delta t = T\).

And again, we will approximate this a bit more by breaking up each time
step into each operator: \[
\ket{\psi(T)} \approx 
e^{-i\tfrac{1}{4}\hat{H}_0\Delta t}e^{-i\tfrac{3}{4}\hat{H}_1\Delta t}
e^{-i\tfrac{1}{2}\hat{H}_0\Delta t}e^{-i\tfrac{3}{2}\hat{H}_1\Delta t}
e^{-i\tfrac{3}{4}\hat{H}_0\Delta t}e^{-i\tfrac{1}{4}\hat{H}_1\Delta t}
e^{-i\hat{H}_0\Delta t}           
\ket{\psi(0)},
\] Here the basic idea is to use a quantum computer to do this time
evolution for us. We will use the fact that Qiskit defines for us the
following gates (or operations):

Single qubit gates (applied to a given qubit \texttt{n}): -
\texttt{h(n)}: \(\hat{U}^H\) - \texttt{rx(θ,n)}:
\(e^{-i\tfrac{\theta}{2} \hat{x}}\) - \texttt{rz(θ,n)}:
\(e^{-i\tfrac{\theta}{2} \hat{z}}\)

Two qubit gates (applied to two given qubits, \texttt{n} and
\texttt{m}): - \texttt{rzz(θ,n,m)}:
\(e^{-i\tfrac{\theta}{2} \hat{z}_n\otimes \hat{z}_m}\)

\subsubsection{Trotterization}\label{trotterization}

One trick we will use is \texttt{Trotterization}. For 2 operators
\(\hat{o}_1\) and \(\hat{o}_2\), we can say the following two things: -
if \(\hat{o}_1\hat{o}_2 = \hat{o}_2\hat{o}_1\), then
\(e^{a\left(\hat{o}_1 + \hat{o}_2\right)} = e^{a\hat{o}_1}e^{a\hat{o}_2}\)
- if \(\hat{o}_1\hat{o}_2 \neq \hat{o}_2\hat{o}_1\) but \(a\ll 1\), then
\(e^{a\left(\hat{o}_1 + \hat{o}_2\right)} \approx e^{a\hat{o}_1}e^{a\hat{o}_2}\)

Where we say that \(\hat{o}_1\) and \(\hat{o}_2\) \texttt{commute} if
\(\hat{o}_1\hat{o}_2 = \hat{o}_2\hat{o}_1\)

\subsection{Initialization}\label{initialization}

From our discussion before, we saw that applying a Hadamard gate to the
\(\ket{0}\) state created the \(\ket{+}\) state, which happened to be
the +1 eigenstate of \(\sigma_x\): - \(\sigma_x\ket{+} = \ket{+}\)

As such, we can initialize our system in the ground state of
\(\hat{H}_0\) by simply applying a Hadamard to each qubit.

    \begin{center}\rule{0.5\linewidth}{0.5pt}\end{center}

For this notebook, we'll use Qiskit, which will need to be installed
into your environment:

\begin{verbatim}
conda install -c conda-forge qiskit qiskit-aer pylatexenc
\end{verbatim}

\begin{center}\rule{0.5\linewidth}{0.5pt}\end{center}

    \begin{tcolorbox}[breakable, size=fbox, boxrule=1pt, pad at break*=1mm,colback=cellbackground, colframe=cellborder]
\prompt{In}{incolor}{2}{\boxspacing}
\begin{Verbatim}[commandchars=\\\{\}]
\PY{k+kn}{import}\PY{+w}{ }\PY{n+nn}{numpy}\PY{+w}{ }\PY{k}{as}\PY{+w}{ }\PY{n+nn}{np}
\PY{k+kn}{import}\PY{+w}{ }\PY{n+nn}{matplotlib}\PY{n+nn}{.}\PY{n+nn}{pyplot}\PY{+w}{ }\PY{k}{as}\PY{+w}{ }\PY{n+nn}{plt}
\PY{k+kn}{from}\PY{+w}{ }\PY{n+nn}{qiskit}\PY{+w}{ }\PY{k+kn}{import} \PY{o}{*}
\PY{k+kn}{from}\PY{+w}{ }\PY{n+nn}{qiskit}\PY{n+nn}{.}\PY{n+nn}{visualization}\PY{+w}{ }\PY{k+kn}{import} \PY{n}{plot\PYZus{}histogram}
\PY{k+kn}{from}\PY{+w}{ }\PY{n+nn}{qiskit}\PY{+w}{ }\PY{k+kn}{import} \PY{n}{QuantumCircuit}

\PY{n}{circ} \PY{o}{=} \PY{n}{QuantumCircuit}\PY{p}{(}\PY{l+m+mi}{2}\PY{p}{,}\PY{l+m+mi}{0}\PY{p}{)}
\PY{n}{circ}\PY{o}{.}\PY{n}{h}\PY{p}{(}\PY{l+m+mi}{0}\PY{p}{)}
\PY{n}{circ}\PY{o}{.}\PY{n}{h}\PY{p}{(}\PY{l+m+mi}{1}\PY{p}{)}
\PY{n}{circ}\PY{o}{.}\PY{n}{draw}\PY{p}{(}\PY{l+s+s1}{\PYZsq{}}\PY{l+s+s1}{mpl}\PY{l+s+s1}{\PYZsq{}}\PY{p}{)}
\end{Verbatim}
\end{tcolorbox}
 
            
\prompt{Out}{outcolor}{2}{}
    
    \begin{center}
    \adjustimage{max size={0.9\linewidth}{0.9\paperheight}}{adiabatic_files/adiabatic_9_0.png}
    \end{center}
    { \hspace*{\fill} \\}
    

    \begin{tcolorbox}[breakable, size=fbox, boxrule=1pt, pad at break*=1mm,colback=cellbackground, colframe=cellborder]
\prompt{In}{incolor}{3}{\boxspacing}
\begin{Verbatim}[commandchars=\\\{\}]
\PY{k+kn}{from}\PY{+w}{ }\PY{n+nn}{qiskit\PYZus{}aer}\PY{n+nn}{.}\PY{n+nn}{primitives}\PY{+w}{ }\PY{k+kn}{import} \PY{n}{SamplerV2} \PY{k}{as} \PY{n}{Sampler}
\PY{n}{sampler} \PY{o}{=} \PY{n}{Sampler}\PY{p}{(}\PY{p}{)}
\PY{n}{circ}\PY{o}{.}\PY{n}{measure\PYZus{}all}\PY{p}{(}\PY{p}{)}
\PY{n}{result} \PY{o}{=} \PY{n}{sampler}\PY{o}{.}\PY{n}{run}\PY{p}{(}\PY{p}{[}\PY{n}{circ}\PY{p}{]}\PY{p}{,} \PY{n}{shots}\PY{o}{=}\PY{l+m+mi}{1000}\PY{p}{)}\PY{o}{.}\PY{n}{result}\PY{p}{(}\PY{p}{)}
\PY{n}{counts} \PY{o}{=} \PY{n}{result}\PY{p}{[}\PY{l+m+mi}{0}\PY{p}{]}\PY{o}{.}\PY{n}{data}\PY{o}{.}\PY{n}{meas}\PY{o}{.}\PY{n}{get\PYZus{}counts}\PY{p}{(}\PY{p}{)} 
\PY{n}{plot\PYZus{}histogram}\PY{p}{(}\PY{n}{counts}\PY{p}{)}
\end{Verbatim}
\end{tcolorbox}
 
            
\prompt{Out}{outcolor}{3}{}
    
    \begin{center}
    \adjustimage{max size={0.9\linewidth}{0.9\paperheight}}{adiabatic_files/adiabatic_10_0.png}
    \end{center}
    { \hspace*{\fill} \\}
    

    \subsection{First time step}\label{first-time-step}

How do we implement \(e^{-i\hat{H}_0\Delta t}\) with gates? Let's make
things more concrete first, and consider a simple 2-site Ising problem.
- \$ \hat{H}\_0 = - \sigma\^{}x\_0 - \sigma\^{}x\_1\$ - \$ \hat{H}\_1 =
J\sigma\textsuperscript{z\_0\sigma}z\_1 + \mu\sigma\^{}z\_0 +
\mu\sigma\^{}z\_1\$ - \(J = 1\) - \(\mu = .1\)

or by simplifying the notation a bit: - \$ \hat{H}\_0 = - \hat{x}\_0 -
\hat{x}\_1\$ - \$ \hat{H}\_1 = J\hat{z}\_0\hat{z}\_1 + \mu\hat{z}\_0 +
\mu\hat{z}\_1\$ - \(J = 1\) - \(\mu = .1\)

Because \(\hat{x}_0\) and \(\hat{x}_1\) act on different qubits, they
commute, meaning that we can break this first time step up into 2
sequential steps exactly: \[ 
e^{i\left(\hat{x}_0 + \hat{x}_1\right)\Delta t} = e^{i\hat{x}_0\Delta t} e^{i\hat{x}_1\Delta t} 
\]

Let's choose a small time step, say \(\Delta t = .1\). Remembering the
fact that \texttt{rx(θ,n)}: \(e^{-i\tfrac{\theta}{2} \hat{x}_n}\), we
can use the following quantum circuit to implement this time step:

    \begin{tcolorbox}[breakable, size=fbox, boxrule=1pt, pad at break*=1mm,colback=cellbackground, colframe=cellborder]
\prompt{In}{incolor}{4}{\boxspacing}
\begin{Verbatim}[commandchars=\\\{\}]
\PY{n}{circ} \PY{o}{=} \PY{n}{QuantumCircuit}\PY{p}{(}\PY{l+m+mi}{2}\PY{p}{,}\PY{l+m+mi}{0}\PY{p}{)}
\PY{n}{circ}\PY{o}{.}\PY{n}{h}\PY{p}{(}\PY{l+m+mi}{0}\PY{p}{)}
\PY{n}{circ}\PY{o}{.}\PY{n}{h}\PY{p}{(}\PY{l+m+mi}{1}\PY{p}{)}
\PY{n}{circ}\PY{o}{.}\PY{n}{rx}\PY{p}{(}\PY{o}{\PYZhy{}}\PY{l+m+mi}{2}\PY{o}{*}\PY{l+m+mf}{.1}\PY{p}{,}\PY{l+m+mi}{1}\PY{p}{)}
\PY{n}{circ}\PY{o}{.}\PY{n}{rx}\PY{p}{(}\PY{o}{\PYZhy{}}\PY{l+m+mi}{2}\PY{o}{*}\PY{l+m+mf}{.1}\PY{p}{,}\PY{l+m+mi}{0}\PY{p}{)}
\PY{n}{circ}\PY{o}{.}\PY{n}{measure\PYZus{}all}\PY{p}{(}\PY{p}{)}
\PY{n}{circ}\PY{o}{.}\PY{n}{draw}\PY{p}{(}\PY{n}{output} \PY{o}{=} \PY{l+s+s1}{\PYZsq{}}\PY{l+s+s1}{mpl}\PY{l+s+s1}{\PYZsq{}}\PY{p}{,} \PY{n}{style}\PY{o}{=}\PY{l+s+s2}{\PYZdq{}}\PY{l+s+s2}{iqp}\PY{l+s+s2}{\PYZdq{}}\PY{p}{)}
\end{Verbatim}
\end{tcolorbox}
 
            
\prompt{Out}{outcolor}{4}{}
    
    \begin{center}
    \adjustimage{max size={0.9\linewidth}{0.9\paperheight}}{adiabatic_files/adiabatic_12_0.png}
    \end{center}
    { \hspace*{\fill} \\}
    

    \begin{tcolorbox}[breakable, size=fbox, boxrule=1pt, pad at break*=1mm,colback=cellbackground, colframe=cellborder]
\prompt{In}{incolor}{5}{\boxspacing}
\begin{Verbatim}[commandchars=\\\{\}]
\PY{n}{nshots} \PY{o}{=} \PY{l+m+mi}{10000}
\PY{n}{plot\PYZus{}histogram}\PY{p}{(}\PY{n}{sampler}\PY{o}{.}\PY{n}{run}\PY{p}{(}\PY{p}{[}\PY{n}{circ}\PY{p}{]}\PY{p}{,} \PY{n}{shots}\PY{o}{=}\PY{n}{nshots}\PY{p}{)}\PY{o}{.}\PY{n}{result}\PY{p}{(}\PY{p}{)}\PY{p}{[}\PY{l+m+mi}{0}\PY{p}{]}\PY{o}{.}\PY{n}{data}\PY{o}{.}\PY{n}{meas}\PY{o}{.}\PY{n}{get\PYZus{}counts}\PY{p}{(}\PY{p}{)}\PY{p}{)}
\end{Verbatim}
\end{tcolorbox}
 
            
\prompt{Out}{outcolor}{5}{}
    
    \begin{center}
    \adjustimage{max size={0.9\linewidth}{0.9\paperheight}}{adiabatic_files/adiabatic_13_0.png}
    \end{center}
    { \hspace*{\fill} \\}
    

    \section{\texorpdfstring{\textbf{Question}:}{Question:}}\label{question}

Explain why the measured probability stayed the same in the above plot:

    After the Hadamards each qubit is in \(\ket{+}\), which is an eigenstate
of \(\hat{x}\) with eigenvalue \(+1\). The gate
\(\mathrm{rx}(\theta) = e^{-i(\theta/2)\hat{x}}\) is just a rotation
about the \(\hat{x}\)-axis, so

\[\mathrm{rx}(\theta)\ket{+} = e^{-i\theta/2}\ket{+},\]

ex) the state is unchanged up to a global phase. Since the measurement
is in the \(\hat{z}\)-basis, \(\ket{+}\) still gives probability
\(\tfrac{1}{2}\) for each outcome, and the two-qubit histogram stays
uniform over the four bitstrings. The \(e^{-i\hat{H}_0\Delta t}\) step
has no effect on the ground state of \(\hat{H}_0\).

    \begin{center}\rule{0.5\linewidth}{0.5pt}\end{center}

\subsection{Second time step}\label{second-time-step}

Now we must implement the following operator:
\(e^{-i\left(.75\hat{H}_0 + .25\hat{H}_1\right)\Delta t}\).

Unfortunately, \(zx \neq xz\) and so we can't exactly trotterize this.
However, because \(\Delta t \ll 1\), we can approximate this in a
product form, and this will become increasingly more accurate as we take
smaller time steps.

\[ \hat{H}_0 = - \hat{x}_0 - \hat{x}_1\]
\[ \hat{H}_1 = J\hat{z}_0\hat{z}_1 + \mu\hat{z}_0 + \mu\hat{z}_1\]

\[
e^{-i\left(- \tfrac{3}{4}\hat{x}_0 - \tfrac{3}{4}\hat{x}_1 + \tfrac{1}{4} J\hat{z}_0\hat{z}_1 + \tfrac{1}{4}\mu\hat{z}_0 + \tfrac{1}{4}\mu\hat{z}_1\right)\Delta t}
\approx
e^{i\tfrac{3}{4}\hat{x}_0\Delta t} 
e^{i\tfrac{3}{4}\hat{x}_1\Delta t}
e^{-i \tfrac{1}{4} J\hat{z}_0\hat{z}_1 \Delta t}
e^{-i \tfrac{1}{4} \mu\hat{z}_0 \Delta t}
e^{-i \tfrac{1}{4} \mu\hat{z}_1 \Delta t}
\]

Notice here that we are using a new gate, \texttt{RZZ}. This is a
2-qubit gate that rotates about the product of two \(\hat{z}\)
operators.

    \begin{tcolorbox}[breakable, size=fbox, boxrule=1pt, pad at break*=1mm,colback=cellbackground, colframe=cellborder]
\prompt{In}{incolor}{6}{\boxspacing}
\begin{Verbatim}[commandchars=\\\{\}]
\PY{n}{circ} \PY{o}{=} \PY{n}{QuantumCircuit}\PY{p}{(}\PY{l+m+mi}{2}\PY{p}{,}\PY{l+m+mi}{0}\PY{p}{)}
\PY{c+c1}{\PYZsh{} initialize}
\PY{n}{circ}\PY{o}{.}\PY{n}{h}\PY{p}{(}\PY{l+m+mi}{0}\PY{p}{)}
\PY{n}{circ}\PY{o}{.}\PY{n}{h}\PY{p}{(}\PY{l+m+mi}{1}\PY{p}{)}

\PY{c+c1}{\PYZsh{} time step 1}
\PY{n}{circ}\PY{o}{.}\PY{n}{rx}\PY{p}{(}\PY{o}{\PYZhy{}}\PY{l+m+mi}{2}\PY{o}{*}\PY{l+m+mf}{.1}\PY{p}{,}\PY{l+m+mi}{1}\PY{p}{)}
\PY{n}{circ}\PY{o}{.}\PY{n}{rx}\PY{p}{(}\PY{o}{\PYZhy{}}\PY{l+m+mi}{2}\PY{o}{*}\PY{l+m+mf}{.1}\PY{p}{,}\PY{l+m+mi}{0}\PY{p}{)}

\PY{c+c1}{\PYZsh{} time step 2}
\PY{n}{circ}\PY{o}{.}\PY{n}{rz}\PY{p}{(}\PY{l+m+mi}{2} \PY{o}{*} \PY{l+m+mf}{.1} \PY{o}{*} \PY{l+m+mf}{.25} \PY{o}{*} \PY{l+m+mf}{.1}\PY{p}{,}\PY{l+m+mi}{1}\PY{p}{)}
\PY{n}{circ}\PY{o}{.}\PY{n}{rz}\PY{p}{(}\PY{l+m+mi}{2} \PY{o}{*} \PY{l+m+mf}{.1} \PY{o}{*} \PY{l+m+mf}{.25} \PY{o}{*} \PY{l+m+mf}{.1}\PY{p}{,}\PY{l+m+mi}{0}\PY{p}{)}
\PY{n}{circ}\PY{o}{.}\PY{n}{rzz}\PY{p}{(}\PY{l+m+mi}{2} \PY{o}{*} \PY{l+m+mf}{1.} \PY{o}{*} \PY{l+m+mf}{.25} \PY{o}{*} \PY{l+m+mf}{.1}\PY{p}{,}\PY{l+m+mi}{0}\PY{p}{,}\PY{l+m+mi}{1}\PY{p}{)}
\PY{n}{circ}\PY{o}{.}\PY{n}{rx}\PY{p}{(}\PY{o}{\PYZhy{}}\PY{l+m+mi}{2} \PY{o}{*} \PY{l+m+mf}{.75} \PY{o}{*} \PY{l+m+mf}{.1}\PY{p}{,}\PY{l+m+mi}{1}\PY{p}{)}
\PY{n}{circ}\PY{o}{.}\PY{n}{rx}\PY{p}{(}\PY{o}{\PYZhy{}}\PY{l+m+mi}{2} \PY{o}{*} \PY{l+m+mf}{.75} \PY{o}{*} \PY{l+m+mf}{.1}\PY{p}{,}\PY{l+m+mi}{0}\PY{p}{)}
\PY{n}{circ}\PY{o}{.}\PY{n}{measure\PYZus{}all}\PY{p}{(}\PY{p}{)}
\PY{n}{circ}\PY{o}{.}\PY{n}{draw}\PY{p}{(}\PY{n}{output} \PY{o}{=} \PY{l+s+s1}{\PYZsq{}}\PY{l+s+s1}{mpl}\PY{l+s+s1}{\PYZsq{}}\PY{p}{)}
\end{Verbatim}
\end{tcolorbox}
 
            
\prompt{Out}{outcolor}{6}{}
    
    \begin{center}
    \adjustimage{max size={0.9\linewidth}{0.9\paperheight}}{adiabatic_files/adiabatic_17_0.png}
    \end{center}
    { \hspace*{\fill} \\}
    

    From here we can recognize the pattern and start to generalize with a
function!

\subsection{Automate the circuit
building}\label{automate-the-circuit-building}

Because each step will have the same gates, but different angles, we can
simply write a function to create the circuit.

    \begin{tcolorbox}[breakable, size=fbox, boxrule=1pt, pad at break*=1mm,colback=cellbackground, colframe=cellborder]
\prompt{In}{incolor}{7}{\boxspacing}
\begin{Verbatim}[commandchars=\\\{\}]
\PY{k}{def}\PY{+w}{ }\PY{n+nf}{form\PYZus{}circuit}\PY{p}{(}\PY{n}{N}\PY{p}{,} \PY{n}{T}\PY{p}{,} \PY{n}{ham}\PY{p}{)}\PY{p}{:}
    \PY{n}{nq} \PY{o}{=} \PY{n}{ham}\PY{o}{.}\PY{n}{N}
    \PY{n}{dt} \PY{o}{=} \PY{n}{T} \PY{o}{/} \PY{n}{N}
    \PY{n}{circ} \PY{o}{=} \PY{n}{QuantumCircuit}\PY{p}{(}\PY{n}{nq}\PY{p}{,} \PY{l+m+mi}{0}\PY{p}{)}

    \PY{c+c1}{\PYZsh{} Initialize in ground state of H\PYZus{}0:  |+\PYZgt{}\PYZca{}n}
    \PY{k}{for} \PY{n}{i} \PY{o+ow}{in} \PY{n+nb}{range}\PY{p}{(}\PY{n}{nq}\PY{p}{)}\PY{p}{:}
        \PY{n}{circ}\PY{o}{.}\PY{n}{h}\PY{p}{(}\PY{n}{i}\PY{p}{)}

    \PY{k}{for} \PY{n}{k} \PY{o+ow}{in} \PY{n+nb}{range}\PY{p}{(}\PY{n}{N}\PY{p}{)}\PY{p}{:}
        \PY{c+c1}{\PYZsh{} weight on H\PYZus{}0}
        \PY{n}{a} \PY{o}{=} \PY{l+m+mf}{1.0} \PY{o}{\PYZhy{}} \PY{n}{k} \PY{o}{/} \PY{n}{N}
        \PY{c+c1}{\PYZsh{} weight on H\PYZus{}1}
        \PY{n}{b} \PY{o}{=} \PY{n}{k} \PY{o}{/} \PY{n}{N}

        \PY{c+c1}{\PYZsh{} e\PYZca{}\PYZob{}\PYZhy{}i dt b H\PYZus{}1\PYZcb{}}
        \PY{k}{if} \PY{n}{b} \PY{o}{\PYZgt{}} \PY{l+m+mf}{0.0}\PY{p}{:}
            \PY{c+c1}{\PYZsh{} local mu z\PYZus{}i:  rz(2 * mu\PYZus{}i * b * dt)}
            \PY{k}{for} \PY{n}{i} \PY{o+ow}{in} \PY{n+nb}{range}\PY{p}{(}\PY{n}{nq}\PY{p}{)}\PY{p}{:}
                \PY{k}{if} \PY{n}{ham}\PY{o}{.}\PY{n}{mu}\PY{p}{[}\PY{n}{i}\PY{p}{]} \PY{o}{!=} \PY{l+m+mf}{0.0}\PY{p}{:}
                    \PY{n}{circ}\PY{o}{.}\PY{n}{rz}\PY{p}{(}\PY{l+m+mf}{2.0} \PY{o}{*} \PY{n}{ham}\PY{o}{.}\PY{n}{mu}\PY{p}{[}\PY{n}{i}\PY{p}{]} \PY{o}{*} \PY{n}{b} \PY{o}{*} \PY{n}{dt}\PY{p}{,} \PY{n}{i}\PY{p}{)}
            \PY{c+c1}{\PYZsh{} J\PYZus{}ij z\PYZus{}i z\PYZus{}j:  rzz(2 * J\PYZus{}ij * b * dt)}
            \PY{k}{for} \PY{p}{(}\PY{n}{i}\PY{p}{,} \PY{n}{j}\PY{p}{)} \PY{o+ow}{in} \PY{n}{ham}\PY{o}{.}\PY{n}{G}\PY{o}{.}\PY{n}{edges}\PY{p}{(}\PY{p}{)}\PY{p}{:}
                \PY{n}{Jij} \PY{o}{=} \PY{n}{ham}\PY{o}{.}\PY{n}{G}\PY{o}{.}\PY{n}{edges}\PY{p}{[}\PY{n}{i}\PY{p}{,} \PY{n}{j}\PY{p}{]}\PY{p}{[}\PY{l+s+s1}{\PYZsq{}}\PY{l+s+s1}{weight}\PY{l+s+s1}{\PYZsq{}}\PY{p}{]}
                \PY{n}{circ}\PY{o}{.}\PY{n}{rzz}\PY{p}{(}\PY{l+m+mf}{2.0} \PY{o}{*} \PY{n}{Jij} \PY{o}{*} \PY{n}{b} \PY{o}{*} \PY{n}{dt}\PY{p}{,} \PY{n}{i}\PY{p}{,} \PY{n}{j}\PY{p}{)}

        \PY{c+c1}{\PYZsh{} e\PYZca{}\PYZob{}\PYZhy{}i dt a H\PYZus{}0\PYZcb{} = e\PYZca{}\PYZob{}+i dt a sum\PYZus{}i x\PYZus{}i\PYZcb{}}
        \PY{c+c1}{\PYZsh{} rx(theta) = e\PYZca{}\PYZob{}\PYZhy{}i (theta/2) x\PYZcb{}, so theta = \PYZhy{}2 * a * dt}
        \PY{k}{if} \PY{n}{a} \PY{o}{\PYZgt{}} \PY{l+m+mf}{0.0}\PY{p}{:}
            \PY{k}{for} \PY{n}{i} \PY{o+ow}{in} \PY{n+nb}{range}\PY{p}{(}\PY{n}{nq}\PY{p}{)}\PY{p}{:}
                \PY{n}{circ}\PY{o}{.}\PY{n}{rx}\PY{p}{(}\PY{o}{\PYZhy{}}\PY{l+m+mf}{2.0} \PY{o}{*} \PY{n}{a} \PY{o}{*} \PY{n}{dt}\PY{p}{,} \PY{n}{i}\PY{p}{)}

    \PY{n}{circ}\PY{o}{.}\PY{n}{measure\PYZus{}all}\PY{p}{(}\PY{p}{)}
    \PY{k}{return} \PY{n}{circ}
\end{Verbatim}
\end{tcolorbox}

    \subsection{Demonstrations}\label{demonstrations}

Try out different numbers of steps and different Hamiltonians. -
\(\hat{H}(t) = (1-f(t))\hat{H}_0 + f(t)\hat{H}_1\)

Let's build a random 1D Ising Hamiltonian:

    \begin{tcolorbox}[breakable, size=fbox, boxrule=1pt, pad at break*=1mm,colback=cellbackground, colframe=cellborder]
\prompt{In}{incolor}{8}{\boxspacing}
\begin{Verbatim}[commandchars=\\\{\}]
\PY{o}{\PYZpc{}}\PY{k}{load\PYZus{}ext} autoreload
\PY{o}{\PYZpc{}}\PY{k}{autoreload} 2

\PY{k+kn}{import}\PY{+w}{ }\PY{n+nn}{networkx}\PY{+w}{ }\PY{k}{as}\PY{+w}{ }\PY{n+nn}{nx}
\PY{k+kn}{import}\PY{+w}{ }\PY{n+nn}{matplotlib}\PY{n+nn}{.}\PY{n+nn}{pyplot}\PY{+w}{ }\PY{k}{as}\PY{+w}{ }\PY{n+nn}{plt}
\PY{k+kn}{import}\PY{+w}{ }\PY{n+nn}{random}
\PY{k+kn}{import}\PY{+w}{ }\PY{n+nn}{montecarlo}

\PY{n}{N} \PY{o}{=} \PY{l+m+mi}{3}
\PY{n}{G} \PY{o}{=} \PY{n}{nx}\PY{o}{.}\PY{n}{Graph}\PY{p}{(}\PY{p}{)}
\PY{n}{G}\PY{o}{.}\PY{n}{add\PYZus{}nodes\PYZus{}from}\PY{p}{(}\PY{p}{[}\PY{n}{i} \PY{k}{for} \PY{n}{i} \PY{o+ow}{in} \PY{n+nb}{range}\PY{p}{(}\PY{n}{N}\PY{p}{)}\PY{p}{]}\PY{p}{)}
\PY{n}{G}\PY{o}{.}\PY{n}{add\PYZus{}edges\PYZus{}from}\PY{p}{(}\PY{p}{[}\PY{p}{(}\PY{n}{i}\PY{p}{,}\PY{p}{(}\PY{n}{i}\PY{o}{+}\PY{l+m+mi}{1}\PY{p}{)}\PY{o}{\PYZpc{}} \PY{n}{G}\PY{o}{.}\PY{n}{number\PYZus{}of\PYZus{}nodes}\PY{p}{(}\PY{p}{)} \PY{p}{)} \PY{k}{for} \PY{n}{i} \PY{o+ow}{in} \PY{n+nb}{range}\PY{p}{(}\PY{n}{N}\PY{p}{)}\PY{p}{]}\PY{p}{)}

\PY{c+c1}{\PYZsh{}  let\PYZsq{}s add some random weights to the edges.}
\PY{n}{G}\PY{o}{.}\PY{n}{edges}\PY{p}{[}\PY{p}{(}\PY{l+m+mi}{0}\PY{p}{,} \PY{l+m+mi}{1}\PY{p}{)}\PY{p}{]}\PY{p}{[}\PY{l+s+s2}{\PYZdq{}}\PY{l+s+s2}{weight}\PY{l+s+s2}{\PYZdq{}}\PY{p}{]} \PY{o}{=} \PY{o}{\PYZhy{}}\PY{l+m+mf}{0.5392900545734278}
\PY{n}{G}\PY{o}{.}\PY{n}{edges}\PY{p}{[}\PY{p}{(}\PY{l+m+mi}{0}\PY{p}{,} \PY{l+m+mi}{2}\PY{p}{)}\PY{p}{]}\PY{p}{[}\PY{l+s+s2}{\PYZdq{}}\PY{l+s+s2}{weight}\PY{l+s+s2}{\PYZdq{}}\PY{p}{]} \PY{o}{=} \PY{l+m+mf}{0.3694282570558236}
\PY{n}{G}\PY{o}{.}\PY{n}{edges}\PY{p}{[}\PY{p}{(}\PY{l+m+mi}{1}\PY{p}{,} \PY{l+m+mi}{2}\PY{p}{)}\PY{p}{]}\PY{p}{[}\PY{l+s+s2}{\PYZdq{}}\PY{l+s+s2}{weight}\PY{l+s+s2}{\PYZdq{}}\PY{p}{]} \PY{o}{=} \PY{o}{\PYZhy{}}\PY{l+m+mf}{0.37530144644718777}

\PY{c+c1}{\PYZsh{} \PYZsh{} alternatively, you can do this using the random package.}
\PY{c+c1}{\PYZsh{} for e in G.edges:}
\PY{c+c1}{\PYZsh{}     G.edges[e][\PYZsq{}weight\PYZsq{}] = (random.random()*2 \PYZhy{} 1)}
\PY{c+c1}{\PYZsh{}     print(e[0], e[1], G.edges[e][\PYZdq{}weight\PYZdq{}])}

\PY{c+c1}{\PYZsh{} Now Draw the graph. }
\PY{n}{plt}\PY{o}{.}\PY{n}{figure}\PY{p}{(}\PY{l+m+mi}{1}\PY{p}{)}
\PY{n}{nx}\PY{o}{.}\PY{n}{draw}\PY{p}{(}\PY{n}{G}\PY{p}{,} \PY{n}{with\PYZus{}labels}\PY{o}{=}\PY{k+kc}{True}\PY{p}{,} \PY{n}{font\PYZus{}weight}\PY{o}{=}\PY{l+s+s1}{\PYZsq{}}\PY{l+s+s1}{bold}\PY{l+s+s1}{\PYZsq{}}\PY{p}{)}
\PY{n}{plt}\PY{o}{.}\PY{n}{show}\PY{p}{(}\PY{p}{)}

\PY{n}{ham} \PY{o}{=} \PY{n}{montecarlo}\PY{o}{.}\PY{n}{IsingHamiltonian}\PY{p}{(}\PY{n}{G}\PY{p}{)}
\PY{c+c1}{\PYZsh{} let\PYZsq{}s add local mu values }
\PY{k}{for} \PY{n}{i} \PY{o+ow}{in} \PY{n+nb}{range}\PY{p}{(}\PY{n}{ham}\PY{o}{.}\PY{n}{N}\PY{p}{)}\PY{p}{:}
    \PY{n}{ham}\PY{o}{.}\PY{n}{mu}\PY{p}{[}\PY{n}{i}\PY{p}{]} \PY{o}{=} \PY{n}{random}\PY{o}{.}\PY{n}{random}\PY{p}{(}\PY{p}{)}
\end{Verbatim}
\end{tcolorbox}

    \begin{center}
    \adjustimage{max size={0.9\linewidth}{0.9\paperheight}}{adiabatic_files/adiabatic_21_0.png}
    \end{center}
    { \hspace*{\fill} \\}
    
    \begin{tcolorbox}[breakable, size=fbox, boxrule=1pt, pad at break*=1mm,colback=cellbackground, colframe=cellborder]
\prompt{In}{incolor}{9}{\boxspacing}
\begin{Verbatim}[commandchars=\\\{\}]
\PY{k+kn}{from}\PY{+w}{ }\PY{n+nn}{qiskit}\PY{n+nn}{.}\PY{n+nn}{transpiler}\PY{n+nn}{.}\PY{n+nn}{preset\PYZus{}passmanagers}\PY{+w}{ }\PY{k+kn}{import} \PY{n}{generate\PYZus{}preset\PYZus{}pass\PYZus{}manager}

\PY{n}{T} \PY{o}{=} \PY{l+m+mi}{20}
\PY{n}{N} \PY{o}{=} \PY{l+m+mi}{100}

\PY{c+c1}{\PYZsh{} This is the function you need to implement}
\PY{n}{circ} \PY{o}{=} \PY{n}{form\PYZus{}circuit}\PY{p}{(}\PY{n}{N}\PY{p}{,} \PY{n}{T}\PY{p}{,} \PY{n}{ham}\PY{p}{)}

\PY{n}{pm} \PY{o}{=} \PY{n}{generate\PYZus{}preset\PYZus{}pass\PYZus{}manager}\PY{p}{(}\PY{n}{optimization\PYZus{}level}\PY{o}{=}\PY{l+m+mi}{3}\PY{p}{,} \PY{n}{basis\PYZus{}gates}\PY{o}{=}\PY{p}{[}\PY{l+s+s2}{\PYZdq{}}\PY{l+s+s2}{rz}\PY{l+s+s2}{\PYZdq{}}\PY{p}{,} \PY{l+s+s2}{\PYZdq{}}\PY{l+s+s2}{rx}\PY{l+s+s2}{\PYZdq{}}\PY{p}{,} \PY{l+s+s2}{\PYZdq{}}\PY{l+s+s2}{cx}\PY{l+s+s2}{\PYZdq{}}\PY{p}{]}\PY{p}{)}
\PY{n}{circ} \PY{o}{=} \PY{n}{pm}\PY{o}{.}\PY{n}{run}\PY{p}{(}\PY{n}{circ}\PY{p}{)}
\PY{n}{circ}\PY{o}{.}\PY{n}{draw}\PY{p}{(}\PY{n}{output}\PY{o}{=}\PY{l+s+s2}{\PYZdq{}}\PY{l+s+s2}{mpl}\PY{l+s+s2}{\PYZdq{}}\PY{p}{,} \PY{n}{fold}\PY{o}{=}\PY{l+m+mi}{120}\PY{p}{)}
\end{Verbatim}
\end{tcolorbox}
 
            
\prompt{Out}{outcolor}{9}{}
    
    \begin{center}
    \adjustimage{max size={0.9\linewidth}{0.9\paperheight}}{adiabatic_files/adiabatic_22_0.png}
    \end{center}
    { \hspace*{\fill} \\}
    

    \begin{tcolorbox}[breakable, size=fbox, boxrule=1pt, pad at break*=1mm,colback=cellbackground, colframe=cellborder]
\prompt{In}{incolor}{10}{\boxspacing}
\begin{Verbatim}[commandchars=\\\{\}]
\PY{n}{nshots} \PY{o}{=} \PY{l+m+mi}{100000}
\PY{n}{plot\PYZus{}histogram}\PY{p}{(}\PY{n}{sampler}\PY{o}{.}\PY{n}{run}\PY{p}{(}\PY{p}{[}\PY{n}{circ}\PY{p}{]}\PY{p}{,} \PY{n}{shots}\PY{o}{=}\PY{n}{nshots}\PY{p}{)}\PY{o}{.}\PY{n}{result}\PY{p}{(}\PY{p}{)}\PY{p}{[}\PY{l+m+mi}{0}\PY{p}{]}\PY{o}{.}\PY{n}{data}\PY{o}{.}\PY{n}{meas}\PY{o}{.}\PY{n}{get\PYZus{}counts}\PY{p}{(}\PY{p}{)}\PY{p}{)}
\end{Verbatim}
\end{tcolorbox}
 
            
\prompt{Out}{outcolor}{10}{}
    
    \begin{center}
    \adjustimage{max size={0.9\linewidth}{0.9\paperheight}}{adiabatic_files/adiabatic_23_0.png}
    \end{center}
    { \hspace*{\fill} \\}
    

    \begin{tcolorbox}[breakable, size=fbox, boxrule=1pt, pad at break*=1mm,colback=cellbackground, colframe=cellborder]
\prompt{In}{incolor}{11}{\boxspacing}
\begin{Verbatim}[commandchars=\\\{\}]
\PY{n}{samples} \PY{o}{=} \PY{n}{sampler}\PY{o}{.}\PY{n}{run}\PY{p}{(}\PY{p}{[}\PY{n}{circ}\PY{p}{]}\PY{p}{,} \PY{n}{shots}\PY{o}{=}\PY{n}{nshots}\PY{p}{)}\PY{o}{.}\PY{n}{result}\PY{p}{(}\PY{p}{)}\PY{p}{[}\PY{l+m+mi}{0}\PY{p}{]}\PY{o}{.}\PY{n}{data}\PY{o}{.}\PY{n}{meas}\PY{o}{.}\PY{n}{get\PYZus{}counts}\PY{p}{(}\PY{p}{)}
\PY{n+nb}{max}\PY{p}{(}\PY{n}{samples}\PY{p}{,} \PY{n}{key}\PY{o}{=}\PY{n}{samples}\PY{o}{.}\PY{n}{get}\PY{p}{)}
\end{Verbatim}
\end{tcolorbox}

            \begin{tcolorbox}[breakable, size=fbox, boxrule=.5pt, pad at break*=1mm, opacityfill=0]
\prompt{Out}{outcolor}{11}{\boxspacing}
\begin{Verbatim}[commandchars=\\\{\}]
'111'
\end{Verbatim}
\end{tcolorbox}
        
    \section{Validation}\label{validation}

How well does this match the exact result?

    \begin{tcolorbox}[breakable, size=fbox, boxrule=1pt, pad at break*=1mm,colback=cellbackground, colframe=cellborder]
\prompt{In}{incolor}{12}{\boxspacing}
\begin{Verbatim}[commandchars=\\\{\}]
\PY{k+kn}{import}\PY{+w}{ }\PY{n+nn}{copy}\PY{+w}{ }\PY{k}{as}\PY{+w}{ }\PY{n+nn}{cp}
\PY{n}{configs} \PY{o}{=} \PY{p}{[}\PY{p}{]}
\PY{n}{energies} \PY{o}{=} \PY{p}{[}\PY{p}{]}
\PY{n}{conf} \PY{o}{=} \PY{n}{montecarlo}\PY{o}{.}\PY{n}{BitString}\PY{p}{(}\PY{n}{ham}\PY{o}{.}\PY{n}{N}\PY{p}{)}
\PY{k}{for} \PY{n}{i} \PY{o+ow}{in} \PY{n+nb}{range}\PY{p}{(}\PY{l+m+mi}{2}\PY{o}{*}\PY{o}{*}\PY{n}{ham}\PY{o}{.}\PY{n}{N}\PY{p}{)}\PY{p}{:}
    \PY{n}{conf}\PY{o}{.}\PY{n}{set\PYZus{}integer\PYZus{}config}\PY{p}{(}\PY{n}{i}\PY{p}{)}
    \PY{n}{configs}\PY{o}{.}\PY{n}{append}\PY{p}{(}\PY{n+nb}{str}\PY{p}{(}\PY{n}{conf}\PY{p}{)}\PY{p}{)}
    \PY{n}{energies}\PY{o}{.}\PY{n}{append}\PY{p}{(}\PY{n}{ham}\PY{o}{.}\PY{n}{energy}\PY{p}{(}\PY{n}{conf}\PY{p}{)}\PY{p}{)}
\PY{n}{plt}\PY{o}{.}\PY{n}{bar}\PY{p}{(}\PY{n}{configs}\PY{p}{,} \PY{n}{energies}\PY{p}{)}
\PY{n}{plt}\PY{o}{.}\PY{n}{xlabel}\PY{p}{(}\PY{l+s+s1}{\PYZsq{}}\PY{l+s+s1}{Configuration}\PY{l+s+s1}{\PYZsq{}}\PY{p}{)}
\PY{n}{plt}\PY{o}{.}\PY{n}{ylabel}\PY{p}{(}\PY{l+s+s1}{\PYZsq{}}\PY{l+s+s1}{Energy}\PY{l+s+s1}{\PYZsq{}}\PY{p}{)}
\end{Verbatim}
\end{tcolorbox}

            \begin{tcolorbox}[breakable, size=fbox, boxrule=.5pt, pad at break*=1mm, opacityfill=0]
\prompt{Out}{outcolor}{12}{\boxspacing}
\begin{Verbatim}[commandchars=\\\{\}]
Text(0, 0.5, 'Energy')
\end{Verbatim}
\end{tcolorbox}
        
    \begin{center}
    \adjustimage{max size={0.9\linewidth}{0.9\paperheight}}{adiabatic_files/adiabatic_26_1.png}
    \end{center}
    { \hspace*{\fill} \\}
    
    \textbf{Note:} Since 0 and 1 have opposite meanings between the
classical and quantum examples (0 was spin down classically, while
\(\ket{1}\) was the negative eigenvalue of \(\hat{z}\)), the bitstrings
should be opposite one another. For example, the classical state
\texttt{0101} corresponds to the \texttt{1010} state. However, Qiskit
prints the bitstrings from right to left, whereas we print ours from
left to write, so that means \texttt{1010} corresponds to \texttt{1010},
and \texttt{101} corresponds to \texttt{010}.

    \section{\texorpdfstring{\textbf{Question}:}{Question:}}\label{question}

What are the probabilities for measuring the correct bitstring using the
following number of time steps (i.e., values of \texttt{n\_steps}): 1.
\texttt{n\_steps} = 3 1. \texttt{n\_steps} = 4 1. \texttt{n\_steps} = 8
1. \texttt{n\_steps} = 10

    \begin{tcolorbox}[breakable, size=fbox, boxrule=1pt, pad at break*=1mm,colback=cellbackground, colframe=cellborder]
\prompt{In}{incolor}{13}{\boxspacing}
\begin{Verbatim}[commandchars=\\\{\}]
\PY{c+c1}{\PYZsh{} Sweep over n\PYZus{}steps and look at the probability of measuring the correct}
\PY{c+c1}{\PYZsh{} ground\PYZhy{}state bitstring.  Hold T fixed and vary N.}
\PY{n}{T} \PY{o}{=} \PY{l+m+mi}{20}
\PY{n}{nshots} \PY{o}{=} \PY{l+m+mi}{100000}

\PY{c+c1}{\PYZsh{} Find the exact classical ground state (the bitstring with lowest energy).}
\PY{n}{conf} \PY{o}{=} \PY{n}{montecarlo}\PY{o}{.}\PY{n}{BitString}\PY{p}{(}\PY{n}{ham}\PY{o}{.}\PY{n}{N}\PY{p}{)}
\PY{n}{best\PYZus{}int}\PY{p}{,} \PY{n}{best\PYZus{}E} \PY{o}{=} \PY{l+m+mi}{0}\PY{p}{,} \PY{n+nb}{float}\PY{p}{(}\PY{l+s+s2}{\PYZdq{}}\PY{l+s+s2}{inf}\PY{l+s+s2}{\PYZdq{}}\PY{p}{)}
\PY{k}{for} \PY{n}{i} \PY{o+ow}{in} \PY{n+nb}{range}\PY{p}{(}\PY{l+m+mi}{2} \PY{o}{*}\PY{o}{*} \PY{n}{ham}\PY{o}{.}\PY{n}{N}\PY{p}{)}\PY{p}{:}
    \PY{n}{conf}\PY{o}{.}\PY{n}{set\PYZus{}integer\PYZus{}config}\PY{p}{(}\PY{n}{i}\PY{p}{)}
    \PY{n}{E} \PY{o}{=} \PY{n}{ham}\PY{o}{.}\PY{n}{energy}\PY{p}{(}\PY{n}{conf}\PY{p}{)}
    \PY{k}{if} \PY{n}{E} \PY{o}{\PYZlt{}} \PY{n}{best\PYZus{}E}\PY{p}{:}
        \PY{n}{best\PYZus{}E}\PY{p}{,} \PY{n}{best\PYZus{}int} \PY{o}{=} \PY{n}{E}\PY{p}{,} \PY{n}{i}
\PY{n}{conf}\PY{o}{.}\PY{n}{set\PYZus{}integer\PYZus{}config}\PY{p}{(}\PY{n}{best\PYZus{}int}\PY{p}{)}
\PY{n}{classical\PYZus{}gs} \PY{o}{=} \PY{n+nb}{str}\PY{p}{(}\PY{n}{conf}\PY{p}{)}

\PY{c+c1}{\PYZsh{} Convert to the Qiskit measurement string.}
\PY{c+c1}{\PYZsh{} Classical convention: 0 = spin\PYZhy{}down, 1 = spin\PYZhy{}up (s\PYZus{}i = 2 b\PYZus{}i \PYZhy{} 1).}
\PY{c+c1}{\PYZsh{} Quantum convention here: |1\PYZgt{} is the \PYZhy{}1 eigenstate of z, so spin\PYZhy{}down \PYZhy{}\PYZgt{} |1\PYZgt{}.}
\PY{c+c1}{\PYZsh{} So we flip every bit, then reverse because Qiskit prints qubits right\PYZhy{}to\PYZhy{}left.}
\PY{n}{quantum\PYZus{}gs} \PY{o}{=} \PY{l+s+s2}{\PYZdq{}}\PY{l+s+s2}{\PYZdq{}}\PY{o}{.}\PY{n}{join}\PY{p}{(}\PY{l+s+s2}{\PYZdq{}}\PY{l+s+s2}{1}\PY{l+s+s2}{\PYZdq{}} \PY{k}{if} \PY{n}{c} \PY{o}{==} \PY{l+s+s2}{\PYZdq{}}\PY{l+s+s2}{0}\PY{l+s+s2}{\PYZdq{}} \PY{k}{else} \PY{l+s+s2}{\PYZdq{}}\PY{l+s+s2}{0}\PY{l+s+s2}{\PYZdq{}} \PY{k}{for} \PY{n}{c} \PY{o+ow}{in} \PY{n}{classical\PYZus{}gs}\PY{p}{)}\PY{p}{[}\PY{p}{:}\PY{p}{:}\PY{o}{\PYZhy{}}\PY{l+m+mi}{1}\PY{p}{]}
\PY{n+nb}{print}\PY{p}{(}\PY{l+s+sa}{f}\PY{l+s+s2}{\PYZdq{}}\PY{l+s+s2}{classical ground state: }\PY{l+s+si}{\PYZob{}}\PY{n}{classical\PYZus{}gs}\PY{l+s+si}{\PYZcb{}}\PY{l+s+s2}{   E = }\PY{l+s+si}{\PYZob{}}\PY{n}{best\PYZus{}E}\PY{l+s+si}{:}\PY{l+s+s2}{.4f}\PY{l+s+si}{\PYZcb{}}\PY{l+s+s2}{\PYZdq{}}\PY{p}{)}
\PY{n+nb}{print}\PY{p}{(}\PY{l+s+sa}{f}\PY{l+s+s2}{\PYZdq{}}\PY{l+s+s2}{corresponding Qiskit bitstring: }\PY{l+s+si}{\PYZob{}}\PY{n}{quantum\PYZus{}gs}\PY{l+s+si}{\PYZcb{}}\PY{l+s+se}{\PYZbs{}n}\PY{l+s+s2}{\PYZdq{}}\PY{p}{)}

\PY{n}{pm} \PY{o}{=} \PY{n}{generate\PYZus{}preset\PYZus{}pass\PYZus{}manager}\PY{p}{(}\PY{n}{optimization\PYZus{}level}\PY{o}{=}\PY{l+m+mi}{3}\PY{p}{,} \PY{n}{basis\PYZus{}gates}\PY{o}{=}\PY{p}{[}\PY{l+s+s2}{\PYZdq{}}\PY{l+s+s2}{rz}\PY{l+s+s2}{\PYZdq{}}\PY{p}{,} \PY{l+s+s2}{\PYZdq{}}\PY{l+s+s2}{rx}\PY{l+s+s2}{\PYZdq{}}\PY{p}{,} \PY{l+s+s2}{\PYZdq{}}\PY{l+s+s2}{cx}\PY{l+s+s2}{\PYZdq{}}\PY{p}{]}\PY{p}{)}
\PY{n}{results} \PY{o}{=} \PY{p}{\PYZob{}}\PY{p}{\PYZcb{}}
\PY{k}{for} \PY{n}{n\PYZus{}steps} \PY{o+ow}{in} \PY{p}{[}\PY{l+m+mi}{3}\PY{p}{,} \PY{l+m+mi}{4}\PY{p}{,} \PY{l+m+mi}{8}\PY{p}{,} \PY{l+m+mi}{10}\PY{p}{]}\PY{p}{:}
    \PY{n}{circ} \PY{o}{=} \PY{n}{form\PYZus{}circuit}\PY{p}{(}\PY{n}{n\PYZus{}steps}\PY{p}{,} \PY{n}{T}\PY{p}{,} \PY{n}{ham}\PY{p}{)}
    \PY{n}{circ\PYZus{}t} \PY{o}{=} \PY{n}{pm}\PY{o}{.}\PY{n}{run}\PY{p}{(}\PY{n}{circ}\PY{p}{)}
    \PY{n}{counts} \PY{o}{=} \PY{n}{sampler}\PY{o}{.}\PY{n}{run}\PY{p}{(}\PY{p}{[}\PY{n}{circ\PYZus{}t}\PY{p}{]}\PY{p}{,} \PY{n}{shots}\PY{o}{=}\PY{n}{nshots}\PY{p}{)}\PY{o}{.}\PY{n}{result}\PY{p}{(}\PY{p}{)}\PY{p}{[}\PY{l+m+mi}{0}\PY{p}{]}\PY{o}{.}\PY{n}{data}\PY{o}{.}\PY{n}{meas}\PY{o}{.}\PY{n}{get\PYZus{}counts}\PY{p}{(}\PY{p}{)}
    \PY{n}{p} \PY{o}{=} \PY{n}{counts}\PY{o}{.}\PY{n}{get}\PY{p}{(}\PY{n}{quantum\PYZus{}gs}\PY{p}{,} \PY{l+m+mi}{0}\PY{p}{)} \PY{o}{/} \PY{n}{nshots}
    \PY{n}{results}\PY{p}{[}\PY{n}{n\PYZus{}steps}\PY{p}{]} \PY{o}{=} \PY{n}{p}
    \PY{n+nb}{print}\PY{p}{(}\PY{l+s+sa}{f}\PY{l+s+s2}{\PYZdq{}}\PY{l+s+s2}{n\PYZus{}steps = }\PY{l+s+si}{\PYZob{}}\PY{n}{n\PYZus{}steps}\PY{l+s+si}{:}\PY{l+s+s2}{\PYZgt{}3d}\PY{l+s+si}{\PYZcb{}}\PY{l+s+s2}{:  P(}\PY{l+s+si}{\PYZob{}}\PY{n}{quantum\PYZus{}gs}\PY{l+s+si}{\PYZcb{}}\PY{l+s+s2}{) = }\PY{l+s+si}{\PYZob{}}\PY{n}{p}\PY{l+s+si}{:}\PY{l+s+s2}{.4f}\PY{l+s+si}{\PYZcb{}}\PY{l+s+s2}{\PYZdq{}}\PY{p}{)}
\end{Verbatim}
\end{tcolorbox}

    \begin{Verbatim}[commandchars=\\\{\}]
classical ground state: 000   E = -1.7507
corresponding Qiskit bitstring: 111

n\_steps =   3:  P(111) = 0.1659
n\_steps =   4:  P(111) = 0.4525
    \end{Verbatim}

    \begin{Verbatim}[commandchars=\\\{\}]
n\_steps =   8:  P(111) = 0.3019
n\_steps =  10:  P(111) = 0.1126
    \end{Verbatim}

    \subsection{\texorpdfstring{\textbf{Answer}}{Answer}}\label{answer}

The exact probabilities depend on the (un-seeded) random \texttt{mu}
draws used to build the Ising Hamiltonian, so the numbers shift between
sessions. For the run executed above (\(T=20\), 100k shots, classical
ground state \texttt{000} ↔ Qiskit string \texttt{111}), we found:

\begin{longtable}[]{@{}rrr@{}}
\toprule\noalign{}
\texttt{n\_steps} & \(\Delta t = T/N\) & \(P(\text{ground state})\) \\
\midrule\noalign{}
\endhead
\bottomrule\noalign{}
\endlastfoot
3 & 6.67 & ≈ 0.17 \\
4 & 5.00 & ≈ 0.45 \\
8 & 2.50 & ≈ 0.30 \\
10 & 2.00 & ≈ 0.11 \\
\end{longtable}

The qualitative point is that the probability is \textbf{not monotonic}
in \texttt{n\_steps} and is far from 1: with \(T\) held fixed,
increasing \(N\) shrinks \(\Delta t\) and reduces the first-order
Trotter splitting error, but at these coarse step sizes the splitting
error is still large enough to produce big oscillations in the success
probability. To approach the adiabatic limit (\(P \to 1\)) you have to
take both \(T\) and \(N\) large --- for instance \(T=20,\ N=100\) gives
\(P \approx 0.99\) in the cells above.


    % Add a bibliography block to the postdoc
    
    
    
\end{document}
